\documentclass[../../main.tex]{subfiles}% 注意这里的文件路径不能用 ./main.tex ,否则用latexmk编译子文件会报错
\graphicspath{{\subfix{./image/}}{\subfix{../../image/}}} % 指定图片引用路径,后续可以直接使用图片文件名
% 如果图片存放在上级目录,则使用 ../../image/ 作为备用路径

% 例如:
% \begin{figure}[H]
% \centering
% \includegraphics[width=0.3\textwidth]{图.png}
% \caption{}
% \label{figure:图}
% \end{figure}
% 注意:上述\label{}一定要放在\caption{}之后,否则引用图片序号会只会显示??.

\begin{document}

\section{存在性结果的证明}\label{section:ch07-存在性证明}

\cref{theorem:ch07-7.3.1}的证明依赖下列结果，它是\cref{theorem:ch06-6.6.1}的对应物。如\cref{theorem:ch06-6.6.1}情形，证明需要若干预备结果。

\begin{theorem}\label{theorem:ch07-7.5.1}
设 $T>0$，$S\triangleq\mathbb{R}\times(0,T)$。设 $u_0\in\mathscr{R}_f^+(\mathbb{R})$ 与 $\phi\in\operatorname{Lip}([0,\infty))$ 满足 $\lim_{y\to\infty}\phi(y)/y=0$（即 (A0)--(A1) 成立，见\cref{subsection:ch07-假设}；不失一般性设 $\phi(0)=0$）。设 $\{u_{0n}\}\subseteq L^1(\mathbb{R})\cap L^\infty(\mathbb{R})$ 满足\eqref{align:ch07-7.44}--\eqref{align:ch07-7.45}，$\{u_n\}$ 为问题 (Pn) 的熵解序列（见\cref{definition:ch07-Pn熵解}，其存在性由\cref{proposition:ch07-7.4.7}给出）。则存在子列 $\{u_{n_k}\}\subseteq\{u_n\}$，具有下列性质：
\begin{enumerate}[(i)]
\item 存在 $u\in L^\infty_{w^*}(0,T;\mathscr{R}_f^+(\mathbb{R}))$ 使
\begin{align}
u_{n_k}\stackrel{*}{\rightharpoonup}u\quad\text{在 }L^\infty_{w^*}(0,T;\mathscr{R}_f(\mathbb{R})); \label{align:ch07-7.88}
\end{align}
\item 与 $\{u_{n_k}\}$ 相伴的 Young 测度序列 $\{\nu_{n_k}\}\subseteq\mathcal{Y}^+(S;\mathbb{R})$ 窄收敛于 Young 测度 $\nu\in\mathcal{Y}^+(S;\mathbb{R})$，且对几乎所有 $(x,t)\in S$ 有
\begin{align}
u_r(x,t)=\int_{[0,\infty)}y\,\mathrm{d}\nu_{(x,t)}(y), \label{align:ch07-7.89}
\end{align}
\begin{align}
\phi(u_r(x,t))=\int_{[0,\infty)}\phi(y)\,\mathrm{d}\nu_{(x,t)}(y), \label{align:ch07-7.90}
\end{align}
其中 $\{\nu_{(x,t)}\}_{(x,t)\in S}$ 为 $\nu$ 的离散分解，$u_r\in L^\infty(0,T;L^1(\mathbb{R}))$ 为 $u$ 关于 Lebesgue 测度的绝对连续部分的密度，$u_s$ 为 $u$ 的奇异部分，即 $u=u_r\,\mathrm{d}x\,\mathrm{d}t+u_s$；
\item 对每个满足
\begin{align}
\lim_{y\to\infty}\frac{U(y)}{y}=:M_U\in\mathbb{R} \label{align:ch07-7.91}
\end{align}
的 $U\in C([0,\infty))$ 与每个 $\zeta\in C([0,T];C_c(\mathbb{R}))$，有
\begin{align}
\lim_{k\to\infty}\iint_S U(u_{n_k})\zeta\,\mathrm{d}x\,\mathrm{d}t
=\iint_S U^*\zeta\,\mathrm{d}x\,\mathrm{d}t+M_U\int_0^T\langle u_s(\cdot,t),\zeta(\cdot,t)\rangle\,\mathrm{d}t, \label{align:ch07-7.92}
\end{align}
其中 $U^*\in L^\infty(0,T;L^1_{\mathrm{loc}}(\mathbb{R}))$ 定义为
\begin{align}
U^*(x,t)\triangleq \int_{[0,\infty)}U(y)\,\mathrm{d}\nu_{(x,t)}(y)
\quad\text{对几乎所有 }(x,t)\in S. \label{align:ch07-7.93}
\end{align}
\end{enumerate}
\end{theorem}

% SUPPLEMENT: parent=theorem:ch07-7.5.1; kind=motivation; location=定理 7.5.1 后的注记; basis=定理 7.5.1 与 (7.71)、(7.93)、(7.108) + 三项陈述各自在定理 7.3.1 证明中的作用、(A0)(A1) 的逐处用途、(7.91) 中极限 $M_U$ 的必要性及其失效边界、$U^*$ 实际属于 $L^\infty(0,T;L^1(\mathbb{R}))$ 的加强
\begin{remark}\phantomsection\label{remark:ch07-7.5.1动机}
\begin{enumerate}[(i)]
\item 引入动机。由\eqref{align:ch07-7.71}中第二个不等式，$\{u_n\}$ 在 $L^\infty(0,T;L^1(\mathbb{R}))$ 中有界，故由 Banach--Alaoglu 定理可抽取弱$^*$收敛子列，其极限只能是测度值的函数 $u$，此即 (i)。但弱$^*$收敛不保持非线性复合：即使 $u_{n_k}\stackrel{*}{\rightharpoonup}u$，一般也没有 $\phi(u_{n_k})\stackrel{*}{\rightharpoonup}\phi(u_r)$，因为 $\phi$ 不是线性的。定理的 (ii) 用 Young 测度 $\nu$ 保存这一丢失的信息：$\nu_{(x,t)}$ 是 $u_{n_k}(x,t)$ 在 $y$ 变量上的“振荡分布”，其重心按\eqref{align:ch07-7.89}给出极限的绝对连续部分 $u_r$，而按\eqref{align:ch07-7.90}可以用 $\nu_{(x,t)}$ 上的积分重构 $\phi(u_r)$，从而绕开对 $\phi(u_{n_k})$ 直接取极限。
\item (iii) 的定量含义。\eqref{align:ch07-7.92}把 $U(u_{n_k})$ 的极限拆成两项：振荡项 $U^*$（由 $\nu$ 决定）与集中项 $M_Uu_s$（由 $u$ 的奇异部分决定，其密度是 $U$ 在无穷远处的斜率 $M_U$）。这正是弱$^*$收敛中“质量逃向 $y=\infty$”的精确刻画：对固定的 $j$，乘子 $f_{j1}$ 把 $U$ 截断在有界范围内，其复合函数的极限只由 $\nu$ 给出；而在集合 $\{u_{n_k}>y_\varepsilon\}$ 上\eqref{align:ch07-7.99}把 $U(u_{n_k})$ 换成 $M_Uu_{n_k}$，后者按 (i) 收敛到 $M_Uu=M_U(u_r\,\mathrm{d}x\,\mathrm{d}t+u_s)$，其中 $u_r$ 部分已并入 $U^*$，剩余部分即 $M_Uu_s$。因此 (iii) 同时是 (i) 与 (ii) 的合成，并给出在\cref{definition:ch07-熵解}与\cref{definition:ch07-Qtau熵解}中足以取极限的形式：取 $U(y)=y$（$M_U=1$）得\eqref{align:ch07-7.138a}，取 $U(y)=\phi(y)$（$M_\phi=0$）得\eqref{align:ch07-7.138b}，取 $U=E_l,F_l$ 则得熵不等式\eqref{align:ch07-7.5}。
\item 假设的作用。(A0) 提供 $u_0\in\mathscr{R}_f^+(\mathbb{R})$，从而使逼近初值 $\{u_{0n}\}$ 可按\eqref{align:ch07-7.44}--\eqref{align:ch07-7.45}选取，并给出\eqref{align:ch07-7.71}中与 $n$ 无关的界 $\|u_0\|_{\mathscr{R}_f(\mathbb{R})}$：这是 (i) 中紧性的唯一来源。(A1) 同时给出 $\phi$ 的 Lipschitz 连续性与次线性增长 $\phi(y)/y\to0$，前者使 $\phi$ 可用于熵条件，后者使 (iii) 对 $U=\phi$ 成立且 $M_\phi=0$；由 (A1) 与 $\phi(0)=0$ 得到的\eqref{align:ch07-7.1}中的常数 $M$ 在\eqref{align:ch07-7.6}、\eqref{align:ch07-7.101}与第 (iii) 项中反复使用。若 $\phi$ 超线性（$\phi(y)/y\to\infty$），则 (iii) 中 $U=\phi$ 不满足\eqref{align:ch07-7.91}，集中项的密度不存在，\eqref{align:ch07-7.92}不能对 $U=\phi$ 使用；这正是本章假定 (A1) 的实质原因，超线性情形须另行处理（见文献~\cite{ref67}）。
\item 条件\eqref{align:ch07-7.91} 的必要性与 $U^*$ 的正则性。\eqref{align:ch07-7.91}要求极限 $M_U$ 存在且有限，而不只是 $U$ 有线性上界：证明中要按 $M_Uu_{n_k}$ 与 $U(u_{n_k})$ 之差作估计，若 $\lim U(y)/y$ 不存在，则该差值在 $\{u_{n_k}>y_\varepsilon\}$ 上没有可用的量级。至于 $U^*$，由\eqref{align:ch07-7.93}、$\nu_{(x,t)}$ 为概率测度及\eqref{align:ch07-7.91}得 $|U^*|\leqslant M(u_b+1)$（$u_b$ 见\eqref{align:ch07-7.95}），故实际上 $U^*\in L^\infty(0,T;L^1(\mathbb{R}))$；陈述中只写 $L^\infty(0,T;L^1_{\mathrm{loc}}(\mathbb{R}))$ 是因为\eqref{align:ch07-7.92}的试验函数 $\zeta$ 的 $x$-支撑紧，两端只需局部可积性即可配对。
\item 记号的一致性。本节只处理 $\Omega=\mathbb{R}$ 的情形，故 (i) 中 $u$ 取值于 $\mathscr{R}_f^+(\mathbb{R})$；这与\eqref{align:ch07-7.108}中的定义一致，后文出现的 $\Omega$ 均按 $\mathbb{R}$ 理解。由于 $u_{n_k}\geqslant0$，其弱$^*$极限非负，故 $u\in L^\infty_{w^*}(0,T;\mathscr{R}_f^+(\mathbb{R}))$ 而不只是 $\mathscr{R}_f(\mathbb{R})$。
\end{enumerate}
\end{remark}

为证\cref{theorem:ch07-7.5.1}，需要下列结果。

\begin{proposition}\label{proposition:ch07-7.5.2}
设 $T>0$，$S\triangleq\mathbb{R}\times(0,T)$，(A0)--(A1) 成立（即 $u_0\in\mathscr{R}_f^+(\mathbb{R})$，$\phi\in\operatorname{Lip}([0,\infty))$，$\lim_{y\to\infty}\phi(y)/y=0$，$\phi(0)=0$，见\cref{subsection:ch07-假设}），$\{u_n\}$ 为问题 (Pn) 的熵解序列（见\cref{definition:ch07-Pn熵解}）。则存在子列 $\{u_{n_k}\}\subseteq\{u_n\}$、Young 测度 $\nu\in\mathcal{Y}^+(S;\mathbb{R})$ 与 $\sigma\in L^\infty_{w^*}(0,T;\mathscr{R}_f^+(\mathbb{R}))$，使
\begin{align}
u_{n_k}\stackrel{*}{\rightharpoonup}u_b+\sigma\quad\text{在 }L^\infty_{w^*}(0,T;\mathscr{R}_f(\mathbb{R})), \label{align:ch07-7.94}
\end{align}
其中 $u_b\in L^\infty(0,T;L^1(\mathbb{R}))$ 定义为
\begin{align}
u_b(x,t)\triangleq \int_{[0,\infty)}y\,\mathrm{d}\nu_{(x,t)}(y)
\quad\text{对几乎所有 }(x,t)\in S, \label{align:ch07-7.95}
\end{align}
而 $\{\nu_{(x,t)}\}_{(x,t)\in S}$ 为 $\nu$ 的离散分解。
\end{proposition}

\begin{proof}
由\eqref{align:ch07-7.71}中第二个不等式，序列 $\{u_n\}$ 在 $L^\infty(0,T;L^1(\mathbb{R}))\subseteq L^\infty_{w^*}(0,T;\mathscr{R}_f(\mathbb{R}))$ 中有界。由\cref{proposition:ch04-4.4.16}与 Banach--Alaoglu 定理，存在 $\mu\in L^\infty_{w^*}(0,T;\mathscr{R}_f^+(\mathbb{R}))$ 使
\begin{align}
u_{n_k}\stackrel{*}{\rightharpoonup}\mu\quad\text{在 }L^\infty_{w^*}(0,T;\mathscr{R}_f(\mathbb{R})). \label{align:ch07-7.96}
\end{align}
另一方面，因 $\{u_{n_k}\}$ 在 $L^1(S)$ 中有界，由\cref{remark:ch05-5.4.4}(ii)，存在 $\{u_{n_k}\}$ 的子列（不重新标号）与 $\sigma\in\mathscr{R}_f^+(S)$ 使
\begin{align}
u_{n_k}\stackrel{*}{\rightharpoonup}u_b+\sigma\quad\text{在 }\mathscr{R}_f(S). \label{align:ch07-7.97}
\end{align}
由\eqref{align:ch07-7.96}--\eqref{align:ch07-7.97}得 $\mu=u_b+\sigma$，结论得证。
\end{proof}

% SUPPLEMENT: parent=proposition:ch07-7.5.2; kind=proof; location=命题 7.5.2 原证明后的补证; basis=原证明以“由 (7.96)--(7.97) 得 $\mu=u_b+\sigma$”一句带过两个弱$^*$极限的识别，未说明 (7.97) 中 Young 测度 $\nu$ 的来源、$\mathscr{R}_f(S)$ 中的收敛如何与 $L^\infty_{w^*}(0,T;\mathscr{R}_f(\mathbb{R}))$ 中的收敛相容，以及 $\sigma$ 的非负性
\begin{proof}
补出原证明中三处未展开的环节：$\nu$ 的来源、两种弱$^*$收敛的相容性、$\sigma$ 的符号。

第一步：Young 测度 $\nu$ 的来源。由\eqref{align:ch07-7.71}，$\{u_{n_k}\}$ 在 $L^\infty(0,T;L^1(\mathbb{R}))$ 中有界且 $u_{n_k}\geqslant0$，故与它相伴的 Young 测度序列 $\{\nu_{n_k}\}\subseteq\mathcal{Y}^+(S;\mathbb{R})$ 有定义且 $\|\nu_{n_k}\|=\|u_{n_k}\|_{L^1(S)}\leqslant T\|u_0\|_{\mathscr{R}_f(\mathbb{R})}$。由\cref{theorem:ch05-Young弱星紧}与\cref{definition:ch05-窄收敛}，可再抽取子列（不重新标号）使 $\nu_{n_k}\stackrel{n}{\rightharpoonup}\nu$ 于 $\mathcal{Y}^+(S;\mathbb{R})$。以下均就这一子列论证；由于命题只断言子列的存在性，再取子列不影响结论。

第二步：\eqref{align:ch07-7.97} 的完整含义。对上述子列应用\cref{remark:ch05-5.4.4}(ii)（取 $X=S$，$\mu$ 为 $S$ 上的 Lebesgue 测度），可能再取子列，存在 $\tilde u\in L^1(S)$ 与 $\sigma\in\mathscr{R}_f(S)$ 使 $u_{n_k}\stackrel{*}{\rightharpoonup}\tilde u+\sigma$ 于 $\mathscr{R}_f(S)$，其中 $\tilde u$ 是 $\nu$ 的重心，即 $\tilde u(x,t)=\int_{[0,\infty)}y\,\mathrm{d}\nu_{(x,t)}(y)=u_b(x,t)$ 对几乎所有 $(x,t)\in S$；这正是\eqref{align:ch07-7.95}。又因 $u_{n_k}\geqslant0$，对每个非负 $\zeta\in C_c(S)$ 有 $\langle u_{n_k},\zeta\rangle\geqslant0$，令 $k\to\infty$ 得 $\langle u_b+\sigma,\zeta\rangle\geqslant0$，故 $\sigma\geqslant0$，即 $\sigma\in\mathscr{R}_f^+(S)$。

第三步：两种收敛的相容性。设 $\zeta\in C([0,T];C_c(\mathbb{R}))$。把 $\zeta$ 看作 $S$ 上的 Borel 可测函数：它在 $x$ 方向有紧支撑、在 $t$ 方向有界，故 $\zeta\in L^1(|\mu|+|\sigma|+u_b\,\mathrm{d}x\,\mathrm{d}t)$，且可被 $C_c(S)$ 中的函数在 $L^1$ 意义下逼近（用 $t$ 方向的截断 $\psi_\delta\in C_c(0,T)$，$\psi_\delta=1$ 于 $(\delta,T-\delta)$，再令 $\delta\to0$ 并用控制收敛定理）。于是由第二步的 $\mathscr{R}_f(S)$ 弱$^*$收敛得
\begin{align*}
\lim_{k\to\infty}\iint_S u_{n_k}\zeta\,\mathrm{d}x\,\mathrm{d}t=\iint_S(u_b+\sigma)\zeta\,\mathrm{d}x\,\mathrm{d}t.
\end{align*}
另一方面，\eqref{align:ch07-7.96}按定义给出同一极限 $\int_0^T\langle\mu(\cdot,t),\zeta(\cdot,t)\rangle\,\mathrm{d}t$，而后者恰是 $\iint_S\mu\,\zeta\,\mathrm{d}x\,\mathrm{d}t$（把 $L^\infty_{w^*}(0,T;\mathscr{R}_f(\mathbb{R}))$ 的元素看作 $S$ 上的测度）。故
\begin{align*}
\iint_S\mu\,\zeta\,\mathrm{d}x\,\mathrm{d}t=\iint_S(u_b+\sigma)\zeta\,\mathrm{d}x\,\mathrm{d}t
\quad\text{对每个 }\zeta\in C([0,T];C_c(\mathbb{R})).
\end{align*}
由于形如 $\zeta(x,t)=h(t)\rho(x)$（$h\in C([0,T])$，$\rho\in C_c(\mathbb{R})$）的函数在 $C([0,T];C_c(\mathbb{R}))$ 中线性张成稠密子集，上式对一切 $\zeta\in C_c(S)$ 成立；测度由它在 $C_c(S)$ 上的值唯一确定，故 $\mu=u_b+\sigma$ 作为 $S$ 上的测度成立，即\eqref{align:ch07-7.94}成立。
\end{proof}

% SUPPLEMENT: parent=proposition:ch07-7.5.2; kind=motivation; location=命题 7.5.2 后的注记; basis=命题 7.5.2 与 (7.94)--(7.95)、(7.71) + 为何把弱$^*$极限拆成“重心部分 + 集中部分”、$\sigma$ 与 Lebesgue 测度的关系、$u_b\in L^\infty(0,T;L^1(\mathbb{R}))$ 而非仅 $L^1(S)$ 的原因，以及 $\sigma$ 不必奇异于 Lebesgue 测度的反例线索
\begin{remark}\phantomsection\label{remark:ch07-7.5.2动机}
\begin{enumerate}[(i)]
\item 引入动机。\cref{proposition:ch07-7.5.2}是“biting 引理”式的分解：$\{u_{n_k}\}$ 在 $L^1(S)$ 中有界，故其弱$^*$极限一般不再是函数。\eqref{align:ch07-7.94}把它拆成两部分：$u_b$ 是 Young 测度 $\nu$ 的重心，属于 $L^\infty(0,T;L^1(\mathbb{R}))$，即“仍可由函数刻画”的振荡部分；$\sigma\in L^\infty_{w^*}(0,T;\mathscr{R}_f^+(\mathbb{R}))$ 是“质量集中”的部分，不再有密度。后续\cref{proposition:ch07-7.5.5}将证明 $\sigma$ 恰是 $u$ 的奇异部分 $u_s$，因此\cref{proposition:ch07-7.5.2}提供了把 (P) 的极限测度拆成绝对连续部分与奇异部分的唯一入口。
\item 各项假设的作用。(A0)--(A1) 只通过\eqref{align:ch07-7.71}的界 $\|u_n\|_{L^1(S)}\leqslant T\|u_0\|_{\mathscr{R}_f(\mathbb{R})}$ 起作用：它同时给出 $L^\infty_{w^*}(0,T;\mathscr{R}_f(\mathbb{R}))$ 中的弱$^*$紧性（\cref{proposition:ch04-4.4.16}与 Banach--Alaoglu 定理）与 $L^1(S)$ 中的有界性（\cref{remark:ch05-5.4.4}(ii) 的前提）。若把界换成只对几乎所有 $t$ 成立而非常数界，则 $\nu$ 的总质量可能无界，$\mathcal{Y}^+$ 中的窄收敛子列不再保证存在。
\item $u_b$ 为何属于 $L^\infty(0,T;L^1(\mathbb{R}))$。\eqref{align:ch07-7.95}只在对几乎所有 $(x,t)$ 定义 $u_b$，但由 Fatou 引理与 $\|\nu_{n_k}\|=\|u_{n_k}\|_{L^1(S)}$ 得
\begin{align*}
\|u_b\|_{L^1(S)}=\iint_S\int_{[0,\infty)}y\,\mathrm{d}\nu_{(x,t)}(y)\,\mathrm{d}x\,\mathrm{d}t\leqslant\liminf_{k\to\infty}\|u_{n_k}\|_{L^1(S)}\leqslant T\|u_0\|_{\mathscr{R}_f(\mathbb{R})},
\end{align*}
且 $\|u_b(\cdot,t)\|_{L^1(\mathbb{R})}\leqslant\|u_0\|_{\mathscr{R}_f(\mathbb{R})}$ 对几乎所有 $t$ 成立，故 $u_b\in L^\infty(0,T;L^1(\mathbb{R}))$。这一点在\cref{proposition:ch07-7.5.5}的证明中给出 $u_b\in L^\infty(0,T;L^1(\mathbb{R}))$ 型估计时被反复使用。
\item 边界。$\sigma$ 一般并不关于 Lebesgue 测度奇异：\cref{remark:ch05-5.4.4}(i) 与\cref{example:ch05-集中非奇异}说明\cref{remark:ch05-5.4.4}(ii) 给出的分解不是 Lebesgue--Radon--Nikodým 分解。\cref{proposition:ch07-7.5.5}将借助熵不等式证明在本问题中 $\sigma$ 确实奇异，从而\eqref{align:ch07-7.108}的分解与 Lebesgue 分解一致；这一步不能由\cref{proposition:ch07-7.5.2}本身得到。
\end{enumerate}
\end{remark}

\begin{proposition}\label{proposition:ch07-7.5.3}
设 $T>0$，$S\triangleq\mathbb{R}\times(0,T)$，(A0)--(A1) 成立（见\cref{subsection:ch07-假设}），$\{u_n\}$ 为问题 (Pn) 的熵解序列。设子列 $\{u_{n_k}\}\subseteq\{u_n\}$、Young 测度 $\nu\in\mathcal{Y}^+(S;\mathbb{R})$ 与 $\sigma\in L^\infty_{w^*}(0,T;\mathscr{R}_f^+(\mathbb{R}))$ 由\cref{proposition:ch07-7.5.2}给出。设 $U\in C([0,\infty))$ 满足\eqref{align:ch07-7.91}，即极限 $M_U\triangleq\lim_{y\to\infty}U(y)/y$ 存在且有限。则对任意 $\zeta\in C([0,T];C_c(\mathbb{R}))$ 有
\begin{align}
\lim_{k\to\infty}\iint_S U(u_{n_k})\zeta\,\mathrm{d}x\,\mathrm{d}t
=\iint_S U^*\zeta\,\mathrm{d}x\,\mathrm{d}t+M_U\int_0^T\langle\sigma(\cdot,t),\zeta(\cdot,t)\rangle\,\mathrm{d}t, \label{align:ch07-7.98}
\end{align}
其中 $U^*\in L^\infty(0,T;L^1_{\mathrm{loc}}(\mathbb{R}))$ 由\eqref{align:ch07-7.93}给出（$\nu$ 为上述 Young 测度）。
\end{proposition}

\begin{proof}
固定任意 $L>0$。由\eqref{align:ch07-7.91}，对任意 $\varepsilon>0$，存在 $y_\varepsilon>0$ 使对所有 $y>y_\varepsilon$ 有
\begin{align}
-\varepsilon y<U(y)-M_U y<\varepsilon y. \label{align:ch07-7.99}
\end{align}
对任意 $j\in\mathbb{N}$、$j\geqslant[y_\varepsilon]+1$，设 $f_{j1},f_{j2}\in C([0,\infty))$ 满足
\begin{align}
\begin{cases}
0\leqslant f_{j1}\leqslant1,\ 0\leqslant f_{j2}\leqslant1,\ f_{j1}+f_{j2}=1\text{ 在 }[0,\infty),\\
\operatorname{supp}f_{j1}\subseteq[0,j+1],\ \operatorname{supp}f_{j2}\subseteq[j,\infty).
\end{cases} 
\label{align:ch07-7.100}
\end{align}
利用\eqref{align:ch07-7.99}与\eqref{align:ch07-7.71}，对固定的 $k\in\mathbb{N}$ 与 $t\in(0,T)$，
\begin{align*}
\begin{aligned}
\int_{I_L}|U(u_{n_k}(x,t))|\,\mathrm{d}x
&=\int_{I_L}f_{j1}(u_{n_k}(x,t))|U(u_{n_k}(x,t))|\,\mathrm{d}x+\int_{I_L}f_{j2}(u_{n_k}(x,t))|U(u_{n_k}(x,t))|\,\mathrm{d}x\\
&\leqslant2L\|U\|_{L^\infty([0,j+1])}+\int_{\{u_{n_k}(\cdot,t)>y_\varepsilon\}}|U(u_{n_k}(x,t))|\,\mathrm{d}x\\
&\leqslant2L\|U\|_{L^\infty([0,j+1])}+(|M_U|+\varepsilon)\int_{I_L}|u_{n_k}(x,t)|\,\mathrm{d}x\\
&\leqslant2L\|U\|_{L^\infty([0,j+1])}+(|M_U|+\varepsilon)\|u_0\|_{\mathscr{R}_f(\mathbb{R})}
\end{aligned}
\end{align*}
（回忆 $I_L=(-L,L)$），故
\begin{align}
\|U(u_{n_k})\|_{L^\infty(0,T;L^1(I_L))}\leqslant2L\|U\|_{L^\infty([0,j+1])}+(|M_U|+\varepsilon)\|u_0\|_{\mathscr{R}_f(\mathbb{R})}. \label{align:ch07-7.101}
\end{align}
由 Banach--Alaoglu 定理，存在 $\{U(u_{n_k})\}$ 的子列（不重新标号）与 $\mu\in L^\infty_{w^*}(0,T;\mathscr{R}_f(I_L))$ 使
\begin{align}
U(u_{n_k})\stackrel{*}{\rightharpoonup}\mu\quad\text{在 }L^\infty_{w^*}(0,T;\mathscr{R}_f(I_L)). \label{align:ch07-7.102}
\end{align}
此外，如\cref{proposition:ch05-5.4.13}证明中那样论证，有
\begin{align}
U(u_{n_k})\stackrel{*}{\rightharpoonup}U^*+M_U\sigma\quad\text{在 }\mathscr{R}_f(R_L). \label{align:ch07-7.103}
\end{align}
由\eqref{align:ch07-7.102}--\eqref{align:ch07-7.103}，对任意 $\zeta\in C([0,T];C_c(I_L))$ 有
\begin{align}
\lim_{k\to\infty}\iint_{R_L}U(u_{n_k})\zeta\,\mathrm{d}x\,\mathrm{d}t
=\iint_{R_L}U^*\zeta\,\mathrm{d}x\,\mathrm{d}t+M_U\int_0^T\langle\sigma(\cdot,t),\zeta(\cdot,t)\rangle\,\mathrm{d}t, \label{align:ch07-7.104}
\end{align}
其中 $U^*$ 由\eqref{align:ch07-7.93}在 $R_L$ 中几乎处处定义。

等式\eqref{align:ch07-7.98}由\eqref{align:ch07-7.104}得到：因对任意 $\zeta\in C([0,T];C_c(\mathbb{R}))$，可取 $L>0$ 充分大使 $\operatorname{supp}\zeta\subseteq((-L,L)\times[0,T])$。此外，由\eqref{align:ch07-7.91}存在 $M>0$ 使 $|U(y)|\leqslant M(y+1)$（$y\geqslant0$），故 $|U^*|\leqslant M(u_b+1)$ 在 $S$ 中几乎处处（见\eqref{align:ch07-7.93}与\eqref{align:ch07-7.95}）。因 $u_b\in L^\infty(0,T;L^1(\mathbb{R}))$，得 $U^*\in L^\infty(0,T;L^1_{\mathrm{loc}}(\mathbb{R}))$。
\end{proof}

% SUPPLEMENT: parent=proposition:ch07-7.5.3; kind=proof; location=命题 7.5.3 原证明中“如命题 5.4.13 证明中那样论证”一句; basis=原书对 (7.103) 只作类比引用；此处核验命题 5.4.13 的全部假设（$X=R_L$、$\mu$ 为 Lebesgue 测度、序列在 $L^1(R_L)$ 中有界、$U$ 在 $\pm\infty$ 处的斜率 $M^\pm_U$），说明 $\sigma_-=0$、$\sigma_+=\sigma$，并由 (7.94) 定出 $U^*$
\begin{proof}
补出\eqref{align:ch07-7.103}的推导。把 $U$ 延拓为 $\mathbb{R}$ 上的连续函数，仍记作 $U$，使 $U(y)=U(0)$ 对 $y\leqslant0$；因 $u_{n_k}\geqslant0$，复合函数 $U(u_{n_k})$ 不变。于是
\begin{align*}
M^+_U\triangleq\lim_{y\to+\infty}\frac{U(y)}{y}=M_U,\qquad
M^-_U\triangleq\lim_{y\to-\infty}\frac{U(y)}{y}=0.
\end{align*}
取 $X=R_L=I_L\times(0,T)$，$\mu$ 为 $X$ 上的 Lebesgue 测度，$\{u_{n_k}\}\subseteq L^1(X)$。由\eqref{align:ch07-7.71}中第二个不等式，$\|u_{n_k}\|_{L^1(X)}\leqslant\|u_0\|_{\mathscr{R}_f(\mathbb{R})}$，故该序列有界，\cref{proposition:ch05-5.4.13}的假设满足。设 $\nu\in\mathcal{Y}^+(X;\mathbb{R})$ 与 $\sigma_\pm\in\mathscr{R}_f^+(X)$ 为该命题中的 Young 测度与两个集中测度；由\cref{proposition:ch05-函数弱星}，可能再取子列，有
\begin{align*}
U(u_{n_k})\stackrel{*}{\rightharpoonup}U^*+M^+_U\sigma_+-M^-_U\sigma_-=U^*+M_U\sigma_+
\quad\text{于 }\mathscr{R}_f(X).
\end{align*}
另一方面，在 $X$ 上应用\cref{proposition:ch07-7.5.2}（其结论\eqref{align:ch07-7.94}限制到 $X$ 后仍成立）得 $u_{n_k}\stackrel{*}{\rightharpoonup}u_b+\sigma$ 于 $\mathscr{R}_f(X)$。对非负序列 $u_{n_k}$，\cref{proposition:ch05-函数弱星}中同一分解给出 $u_{n_k}\stackrel{*}{\rightharpoonup}u^*+\sigma_+-\sigma_-$，其中 $u^*=\int y\,\mathrm{d}\nu$；与前一式比较（测度由其在 $C_c(X)$ 上的值唯一确定）得
\begin{align*}
u_b+\sigma=u^*+\sigma_+-\sigma_-.
\end{align*}
因 $u_{n_k}\geqslant0$，其 Young 测度 $\nu$ 支撑于 $[0,\infty)$，从而没有质量逃向 $-\infty$，即 $\sigma_-=0$；又 $u^*=\int_{[0,\infty)}y\,\mathrm{d}\nu_{(x,t)}(y)=u_b$ 于 $X$ 中几乎处处（\eqref{align:ch07-7.95}）。故 $\sigma_+=\sigma$，\eqref{align:ch07-7.103}成立。将\eqref{align:ch07-7.103}与\eqref{align:ch07-7.102}比较即得 $\mu=U^*+M_U\sigma$（作为 $X$ 上的测度），再由 $\zeta\in C([0,T];C_c(I_L))$ 与\eqref{align:ch07-7.102}的配对即得\eqref{align:ch07-7.104}。
\end{proof}

% SUPPLEMENT: parent=proposition:ch07-7.5.3; kind=motivation; location=命题 7.5.3 后的注记; basis=命题 7.5.3 与 (7.98)、(7.101)、(7.103) + 为何此处出现 $\sigma$ 而不是 $u_s$（逻辑顺序）、(7.99) 的截断策略、$\varepsilon$ 与 $j$ 的依赖次序、$L$ 的局部化作用，以及 (7.91) 不可去掉的边界
\begin{remark}\phantomsection\label{remark:ch07-7.5.3动机}
\begin{enumerate}[(i)]
\item 引入动机。\cref{proposition:ch07-7.5.3}是\cref{theorem:ch07-7.5.1}(iii)的原型：它把\eqref{align:ch07-7.92}右端的 $u_s$ 暂换为尚未识别的集中测度 $\sigma$。这样做的原因是逻辑顺序：$u_s=\sigma$ 的证明（\cref{proposition:ch07-7.5.5}）本身要用到熵不等式\eqref{align:ch07-7.77}，而\eqref{align:ch07-7.77}取极限又需要\eqref{align:ch07-7.98}；故必须先在 $\sigma$ 的层面上建立\eqref{align:ch07-7.98}，再识别 $\sigma$。若直接写 $u_s$，论证将出现循环。
\item 截断策略。(7.99) 把 $U$ 在无穷远处线性化为 $M_Uy$，(7.100) 把 $y$ 轴分成 $[0,j+1]$ 与 $[j,\infty)$ 两部分：在 $[0,j+1]$ 上 $U$ 有界，故 $|U(u_{n_k})|$ 可直接由 $2L\|U\|_{L^\infty([0,j+1])}$ 控制；在 $\{u_{n_k}>y_\varepsilon\}$ 上则用 (7.99) 换成 $(|M_U|+\varepsilon)|u_{n_k}|$。两处都只需\eqref{align:ch07-7.71}的界，与 $k$ 无关，这正是\eqref{align:ch07-7.101}中常数不含 $k$ 的原因，也是 Banach--Alaoglu 定理可用的前提。
\item 量词的次序。$j$ 依赖 $\varepsilon$（要求 $j\geqslant[y_\varepsilon]+1$，使 $\operatorname{supp}f_{j2}\subseteq[y_\varepsilon,\infty)$），而\eqref{align:ch07-7.101}对每个固定的 $j$ 成立；取子列时 $j$ 与 $\varepsilon$ 固定，最后才让 $\varepsilon\downarrow0$。$\varepsilon$ 不能取为 $0$：\eqref{align:ch07-7.99}的严格不等式只在 $y>y_\varepsilon$ 上成立，而 $U(y)-M_Uy$ 在有限区间上一般不为零。
\item 局部化的作用。\eqref{align:ch07-7.101}--\eqref{align:ch07-7.104}只在 $I_L$ 上陈述，原因是 $\int_{\mathbb{R}}|u_{n_k}|\,\mathrm{d}x$ 虽一致有界，但\eqref{align:ch07-7.101}右端的界含 $2L\|U\|_{L^\infty([0,j+1])}$，它随 $L$ 增长；因此只能先在 $R_L$ 上取弱$^*$极限，再借 $\zeta$ 的紧支撑把 $L$ 固定下来。\eqref{align:ch07-7.98}中的 $U^*$ 只需在 $\operatorname{supp}\zeta$ 上几乎处处有定义，故 $L^\infty(0,T;L^1_{\mathrm{loc}}(\mathbb{R}))$ 足够。
\item 边界。\eqref{align:ch07-7.91}不可去掉：若 $U$ 在 $[0,\infty)$ 上只有线性上界而 $\lim_{y\to\infty}U(y)/y$ 不存在，则\eqref{align:ch07-7.99}不成立，(7.101) 中的常数无法与 $k$ 无关。此外 $\sigma$ 非负性在\eqref{align:ch07-7.104}中并未使用，但它是\cref{proposition:ch07-7.5.5}中把 $\sigma$ 与 $u_{0s}$ 比较的前提，故在\cref{proposition:ch07-7.5.2}中一并列出。
\end{enumerate}
\end{remark}

为继续，需要下列引理，其证明在本节末给出。

\begin{lemma}\label{lemma:ch07-7.5.4}
设 $T>0$，$S\triangleq\mathbb{R}\times(0,T)$，(A0)--(A1) 成立（见\cref{subsection:ch07-假设}），$U\in C([0,\infty))$ 满足\eqref{align:ch07-7.91}（即 $M_U\triangleq\lim_{y\to\infty}U(y)/y$ 存在且有限）。设子列 $\{u_{n_k}\}\subseteq\{u_n\}$、$\sigma\in L^\infty_{w^*}(0,T;\mathscr{R}_f^+(\mathbb{R}))$ 与 $U^*\in L^\infty(0,T;L^1_{\mathrm{loc}}(\mathbb{R}))$ 如\cref{proposition:ch07-7.5.3}，其中 $\{u_n\}$ 为问题 (Pn) 的熵解序列。则对每个 $\zeta\in C([0,T];C_c(\mathbb{R}))$ 有
\begin{align}
\lim_{k\to\infty}\int_0^T\left|\int_{\mathbb{R}}U(u_{n_k}(x,t))\zeta(x,t)\,\mathrm{d}x
-\int_{\mathbb{R}}U^*(x,t)\zeta(x,t)\,\mathrm{d}x
-M_U\langle\sigma(\cdot,t),\zeta(\cdot,t)\rangle\right|\mathrm{d}t=0. \label{align:ch07-7.105}
\end{align}
\end{lemma}

% SUPPLEMENT: parent=lemma:ch07-7.5.4; kind=motivation; location=引理 7.5.4 后的注记; basis=引理 7.5.4 与 (7.98)、(7.105)、(7.107)、(7.111) + 为何要把 (7.98) 的积分收敛加强为 $L^1(0,T)$ 中的收敛、绝对值的用途、$U$ 的正则性只需连续而无需 $C^2$ 的原因，以及只能得到几乎处处的 $t$ 而不能得到逐点收敛的边界
\begin{remark}\phantomsection\label{remark:ch07-7.5.4动机}
\begin{enumerate}[(i)]
\item 引入动机。\cref{proposition:ch07-7.5.3}只给出\eqref{align:ch07-7.98}：对每个固定的试验函数 $\zeta$，数值序列 $\iint_SU(u_{n_k})\zeta\,\mathrm{d}x\,\mathrm{d}t$ 收敛。但在\cref{proposition:ch07-7.5.5}的证明中需要在某个时刻 $\tau$ 处理 $\int_{\mathbb{R}}U_j(u_{n_k})(x,\tau)\zeta(x,\tau)\,\mathrm{d}x$，即需要关于 $t$ 的逐点信息。\eqref{align:ch07-7.105}把收敛加强为“$\mathrm{d}t$ 积分意义下的收敛”，从而使“抽取子列使对几乎所有 $t$ 成立”成为可能：此即\cref{remark:ch07-7.5.1}中的\eqref{align:ch07-7.107}。换言之，\cref{lemma:ch07-7.5.4}是把弱$^*$收敛的信息从 $S$ 上转移到每个时刻切片上的桥梁。
\item 绝对值的用途。\eqref{align:ch07-7.105}中被积函数取绝对值后再对 $t$ 积分，因此该式对 $\zeta$ 的依赖是“一致的”：只要 $\|\zeta\|_{L^\infty}$ 小，左端的上极限就小。这一点在\cref{remark:ch07-7.5.1}的稠密性论证中被用来从可数稠密集 $\mathcal{D}\subseteq C_c(\mathbb{R})$ 过渡到全部 $\rho\in C_c(\mathbb{R})$：对 $\rho-\rho_i$ 的贡献可用与 $k$ 无关的常数乘 $\|\rho-\rho_i\|_{L^\infty}$ 控制。若\eqref{align:ch07-7.105}不取绝对值，这一过渡无法进行。
\item 假设的作用。引理只要求 $U\in C([0,\infty))$ 满足\eqref{align:ch07-7.91}，不要求 $U\in C^2$ 或 $U'$ 有界：证明分三步处理，先对 $U\in C^2$ 且 $U'\in L^\infty$ 的 $U$（此时可用\cref{lemma:ch07-7.4.6}的 $BV$ 估计），再对 $U\in C_b([0,\infty))$（磨光），最后对一般的 $U$（用\eqref{align:ch07-7.100}的分解 $U=U_{j1}+U_{j2}$，其中 $U_{j2}$ 在无穷远处近似为 $M_Uy$）。\eqref{align:ch07-7.91}在这里的作用仍是给出线性化斜率 $M_U$；当 $U$ 有界时 $M_U=0$，\eqref{align:ch07-7.105}退化为不含 $\sigma$ 的形式\eqref{align:ch07-7.106a}。
\item 边界。\eqref{align:ch07-7.105}不能改进为“对每个 $t\in(0,T)$ 成立”：由\eqref{align:ch07-7.105}只能得到存在零测集 $N$ 使 $t\notin N$ 时\eqref{align:ch07-7.107}成立，因为 $\{g_k\}$ 在 $L^1(0,T)$ 中收敛只蕴含几乎处处收敛的子列。这一限制在\cref{proposition:ch07-7.5.5}中是容许的：那里的结论\eqref{align:ch07-7.109a}本来就只要求几乎处处成立。
\end{enumerate}
\end{remark}

\begin{remark}\phantomsection\label{remark:ch07-7.5.1}
取 $\zeta(x,t)=\rho(x)$、$\rho\in C_c(\mathbb{R})$，由\eqref{align:ch07-7.105}得
\begin{align}
\lim_{k\to\infty}\int_0^T\left|\int_{\mathbb{R}}U(u_{n_k}(x,t))\rho(x)\,\mathrm{d}x-U_\rho^*(t)-M_U\langle\sigma(\cdot,t),\rho\rangle\right|\mathrm{d}t=0, \label{align:ch07-7.106a}
\end{align}
其中
\begin{align}
U_\rho^*(t)\triangleq \int_{\mathbb{R}}U^*(x,t)\rho(x)\,\mathrm{d}x. \label{align:ch07-7.106b}
\end{align}
由\eqref{align:ch07-7.106a}与 $C_c(\mathbb{R})$ 的可分性，对每个 $L>0$ 与如\cref{lemma:ch07-7.5.4}的 $U$，存在零测集 $N\subseteq(0,T)$ 与 $\{u_{n_k}\}$ 的子列（不重新标号），使对所有 $t\in(0,T)\setminus N$ 有
\begin{align}
U(u_{n_k})(\cdot,t)\stackrel{*}{\rightharpoonup}U^*(\cdot,t)+M_U\sigma(\cdot,t)\quad\text{在 }\mathscr{R}_f(I_L). \label{align:ch07-7.107}
\end{align}
\end{remark}

% SUPPLEMENT: parent=remark:ch07-7.5.1; kind=proof; location=注 7.5.1 中“由 (7.106a) 与 $C_c(\mathbb{R})$ 的可分性”一句; basis=原书未写出由 $L^1(0,T)$ 中的收敛到几乎处处 $t$ 的对角抽取、可数稠密集的选取、以及借 (7.101) 与 $\|\sigma(\cdot,t)\|\leqslant\|\sigma\|$ 从稠密集过渡到全部 $\rho\in C_c(\mathbb{R})$ 的论证
\begin{proof}
补出\eqref{align:ch07-7.107}的推导，分四步。

第一步：可数稠密集。$C_c(\mathbb{R})$ 按 $\|\cdot\|_{L^\infty}$ 可分：取定 $\theta\in C_c^\infty(\mathbb{R})$ 使 $\theta=1$ 于 $[-1,1]$、$\operatorname{supp}\theta\subseteq[-2,2]$，并置
\begin{align*}
\mathcal{D}\triangleq\Big\{\sum_{i=1}^{m}q_i\theta(\cdot/l)\,\rho_i\ \Big|\ m\in\mathbb{N},\ l\in\mathbb{N},\ q_i\in\mathbb{Q},\ \rho_i\in\mathcal{P}\Big\},
\end{align*}
其中 $\mathcal{P}$ 为有理系数多项式之集；则 $\mathcal{D}\subseteq C_c(\mathbb{R})$ 可数且在 $(C_c(\mathbb{R}),\|\cdot\|_{L^\infty})$ 中稠密。

第二步：对固定 $\rho\in C_c(\mathbb{R})$ 抽取子列。置
\begin{align*}
g_k^\rho(t)\triangleq\int_{\mathbb{R}}U(u_{n_k}(x,t))\rho(x)\,\mathrm{d}x-U_\rho^*(t)-M_U\langle\sigma(\cdot,t),\rho\rangle
\quad(t\in(0,T)).
\end{align*}
由\eqref{align:ch07-7.106a}，$\|g_k^\rho\|_{L^1(0,T)}\to0$；故由 Chebyshev 不等式，存在子列 $\{k_m\}$ 使 $g_{k_m}^\rho(t)\to0$ 对几乎所有 $t\in(0,T)$。

第三步：对角抽取。把 $\mathcal{D}$ 排成 $\{\rho_1,\rho_2,\cdots\}$。由第二步与\eqref{align:ch07-7.106a}，可归纳地选取 $k_1<k_2<\cdots$ 使
\begin{align*}
\|g_{k_m}^{\rho_i}\|_{L^1(0,T)}\leqslant\frac1m\quad\text{对 }i=1,\cdots,m.
\end{align*}
置 $N\triangleq\bigcup_{i\in\mathbb{N}}N_i$，其中 $N_i\subseteq(0,T)$ 为零测集使 $g_{k_m}^{\rho_i}(t)\to0$（$m\to\infty$）对所有 $t\notin N_i$ 成立；$N$ 是可数并，故为零测集。于是对每个 $t\in(0,T)\setminus N$ 与每个 $i\in\mathbb{N}$，
\begin{align*}
\lim_{m\to\infty}g_{k_m}^{\rho_i}(t)=0.
\end{align*}

第四步：由稠密性过渡到全部 $\rho$。固定 $t\in(0,T)\setminus N$。对任意 $\rho\in C_c(\mathbb{R})$ 与 $\eta>0$，取 $\rho_i\in\mathcal{D}$ 使 $\|\rho-\rho_i\|_{L^\infty}\leqslant\eta$。由\eqref{align:ch07-7.101}与 $\|U^*(\cdot,t)\|_{L^1(\mathbb{R})}\leqslant\|U^*\|_{L^\infty(0,T;L^1(\mathbb{R}))}$、$\|\sigma(\cdot,t)\|_{\mathscr{R}_f(\mathbb{R})}\leqslant\|\sigma\|_{L^\infty_{w^*}(0,T;\mathscr{R}_f(\mathbb{R}))}$，存在与 $k$ 无关的常数 $C_t$（依赖 $t$、$U$、$u_0$ 与 $\sigma$）使
\begin{align*}
|g_k^\rho(t)-g_k^{\rho_i}(t)|\leqslant C_t\|\rho-\rho_i\|_{L^\infty}\leqslant C_t\eta.
\end{align*}
故 $\limsup_{m\to\infty}|g_{k_m}^\rho(t)|\leqslant C_t\eta$；令 $\eta\downarrow0$ 得 $g_{k_m}^\rho(t)=0$ 的极限成立。

第五步：翻译为弱$^*$收敛。设 $\zeta\in C_c(I_L)$。由第四步，对每个 $t\in(0,T)\setminus N$ 与每个 $\rho\in C_c(\mathbb{R})$，
\begin{align*}
\int_{\mathbb{R}}U(u_{n_k}(x,t))\rho(x)\,\mathrm{d}x\to\int_{\mathbb{R}}U^*(x,t)\rho(x)\,\mathrm{d}x+M_U\langle\sigma(\cdot,t),\rho\rangle.
\end{align*}
由于 $\zeta\in C_c(I_L)\subseteq C_c(\mathbb{R})$，且 $U^*(\cdot,t)+M_U\sigma(\cdot,t)\in L^1(I_L)\subseteq\mathscr{R}_f(I_L)$（$U^*(\cdot,t)\in L^1$ 对几乎所有 $t$ 成立，$\sigma(\cdot,t)\in\mathscr{R}_f$ 对几乎所有 $t$ 成立），上式即\eqref{align:ch07-7.107}。
\end{proof}

% SUPPLEMENT: parent=proposition:ch07-divcurl; kind=motivation; 该对象的动机与条件作用见紧随其后的注记
\begin{proposition}\label{proposition:ch07-7.5.5}
设 $T>0$，$S\triangleq\mathbb{R}\times(0,T)$，(A0)--(A1) 成立（见\cref{subsection:ch07-假设}），$\{u_n\}$ 为问题 (Pn) 的熵解序列。设子列 $\{u_{n_k}\}\subseteq\{u_n\}$ 与 $\sigma\in L^\infty_{w^*}(0,T;\mathscr{R}_f^+(\mathbb{R}))$ 由\cref{proposition:ch07-7.5.2}给出，并按
\begin{align}
u\triangleq u_b+\sigma \label{align:ch07-7.108}
\end{align}
定义 $u\in L^\infty_{w^*}(0,T;\mathscr{R}_f^+(\mathbb{R}))$，其中 $u_b$ 由\eqref{align:ch07-7.95}给出。则
\begin{align}
u_r=u_b\quad\text{在 }S\text{ 中几乎处处}, \label{align:ch07-7.109a}
\end{align}
\begin{align}
u_s=\sigma\quad\text{在 }L^\infty_{w^*}(0,T;\mathscr{R}_f(\mathbb{R})). \label{align:ch07-7.109b}
\end{align}
\end{proposition}

\begin{proof}
对任意非负 $\zeta\in C^1([0,T];C^1_c(\mathbb{R}))$、$\zeta(\cdot,T)=0$，将\eqref{align:ch07-7.77}与\eqref{align:ch07-7.78}相加，得
\begin{align*}
\iint_S\{U_j(u_{n_k})\partial_t\zeta+\Lambda_{U_j}(u_{n_k})\partial_x\zeta\}\,\mathrm{d}x\,\mathrm{d}t
\geqslant-\int_{\mathbb{R}}U_j(u_{0n_k})(x)\zeta(x,0)\,\mathrm{d}x,
\end{align*}
其中对所有 $j\in\mathbb{N}$，
\begin{align}
U_j(y)\triangleq (y-j)\chi_{[j,\infty)}(y),\quad
\Lambda_{U_j}(y)\triangleq \int_0^y U_j'(z)\phi'(z)\,\mathrm{d}z=[\phi(y)-\phi(j)]\chi_{[j,\infty)}(y). \label{align:ch07-7.110}
\end{align}
由标准逼近论证，对上述每个 $\zeta$ 与几乎所有 $\tau\in(0,T)$ 有
\begin{align}
\begin{aligned}
\int_{\mathbb{R}}&U_j(u_{n_k})(x,\tau)\zeta(x,\tau)\,\mathrm{d}x-\int_{\mathbb{R}}U_j(u_{0n_k})(x)\zeta(x,0)\,\mathrm{d}x\\
&\leqslant\iint_{S_\tau}\{U_j(u_{n_k})\partial_t\zeta+\Lambda_{U_j}(u_{n_k})\partial_x\zeta\}\,\mathrm{d}x\,\mathrm{d}t,
\end{aligned} 
\label{align:ch07-7.111}
\end{align}
其中 $S_\tau\triangleq \mathbb{R}\times(0,\tau)$。因 $M_{U_j}=1$、$M_{\Lambda_{U_j}}=0$，由\eqref{align:ch07-7.105}得
\begin{align}
\lim_{k\to\infty}\int_0^\tau\left|\int_{\mathbb{R}}[U_j(u_{n_k})\partial_t\zeta](x,t)\,\mathrm{d}x
-\int_{\mathbb{R}}[U_j^*\partial_t\zeta](x,t)\,\mathrm{d}x
-\langle\sigma(\cdot,t),\partial_t\zeta(\cdot,t)\rangle\right|\mathrm{d}t=0, \label{align:ch07-7.112a}
\end{align}
\begin{align}
\lim_{k\to\infty}\int_0^\tau\left|\int_{\mathbb{R}}[\Lambda_{U_j}(u_{n_k})\partial_x\zeta](x,t)\,\mathrm{d}x
-\int_{\mathbb{R}}[\Lambda_{U_j}^*\partial_x\zeta](x,t)\,\mathrm{d}x\right|\mathrm{d}t=0, \label{align:ch07-7.112b}
\end{align}
其中
\begin{align*}
U_j^*\triangleq \int_{[0,\infty)}U_j(y)\,\mathrm{d}\nu(y),\qquad
\Lambda_{U_j}^*\triangleq \int_{[0,\infty)}\Lambda_{U_j}(y)\,\mathrm{d}\nu(y).
\end{align*}
由\eqref{align:ch07-7.112a}得
\begin{align}
\begin{aligned}
\lim_{k\to\infty}\iint_{S_\tau}&\{U_j(u_{n_k})\partial_t\zeta+\Lambda_{U_j}(u_{n_k})\partial_x\zeta\}\,\mathrm{d}x\,\mathrm{d}t\\
&=\iint_{S_\tau}\{U_j^*\partial_t\zeta+\Lambda_{U_j}^*\partial_x\zeta\}\,\mathrm{d}x\,\mathrm{d}t
+\int_0^\tau\langle\sigma(\cdot,t),\partial_t\zeta(\cdot,t)\rangle\,\mathrm{d}t.
\end{aligned} 
\label{align:ch07-7.113}
\end{align}
另一方面，由\eqref{align:ch07-7.107}与对角论证，存在零测集 $N\subseteq(0,T)$ 与 $\{u_{n_k}\}$ 的子列（不重新标号），使对所有 $\tau\in(0,T)\setminus N$ 与 $j\in\mathbb{N}$，
\begin{align}
\lim_{k\to\infty}\int_{\mathbb{R}}U_j(u_{n_k})(x,\tau)\zeta(x,\tau)\,\mathrm{d}x
=\int_{\mathbb{R}}U_j^*(x,\tau)\zeta(x,\tau)\,\mathrm{d}x+\langle\sigma(\cdot,\tau),\zeta(x,\tau)\rangle. \label{align:ch07-7.114}
\end{align}
此外，序列 $\{U_j(u_{0n_k})-u_{0n_k}\}$ 在 $L^\infty(\mathbb{R})$ 中有界且
\begin{align*}
\lim_{k\to\infty}(U_j(u_{0n_k})-u_{0n_k})=U_j(u_{0r})-u_{0r}\quad\text{在 }\mathbb{R}\text{ 中几乎处处}
\end{align*}
（见\eqref{align:ch07-7.45}），故由\eqref{align:ch07-7.45}得
\begin{align}
\lim_{k\to\infty}\int_{\mathbb{R}}U_j(u_{0n_k})(x)\zeta(x,0)\,\mathrm{d}x
=\int_{\mathbb{R}}U_j(u_{0r})(x)\zeta(x,0)\,\mathrm{d}x+\langle u_{0s},\zeta(\cdot,0)\rangle. \label{align:ch07-7.115}
\end{align}
在\eqref{align:ch07-7.111}中令 $k\to\infty$，由\eqref{align:ch07-7.113}--\eqref{align:ch07-7.115}，对所有 $\tau\in(0,T)\setminus N$ 与 $j\in\mathbb{N}$，得
\begin{align}
\begin{aligned}
\int_{\mathbb{R}}U_j^*(x,\tau)\zeta(x,\tau)\,\mathrm{d}x+\langle\sigma(\cdot,\tau),\zeta(\cdot,\tau)\rangle
\leqslant{}&\iint_{S_\tau}\{U_j^*\partial_t\zeta+\Lambda_{U_j}^*\partial_x\zeta\}\,\mathrm{d}x\,\mathrm{d}t\\
&+\int_0^\tau\langle\sigma(\cdot,t),\partial_t\zeta(\cdot,t)\rangle\,\mathrm{d}t
+\int_{\mathbb{R}}U_j(u_{0r})(x)\zeta(x,0)\,\mathrm{d}x+\langle u_{0s},\zeta(\cdot,0)\rangle.
\end{aligned} 
\label{align:ch07-7.116}
\end{align}
因对所有 $y\geqslant0$ 有 $0\leqslant U_j(y)\leqslant y\chi_{[j,\infty)}(y)$ 与 $|\Lambda_{U_j}(y)|\leqslant My\chi_{[j,\infty)}(y)$（见\eqref{align:ch07-7.1}），故 $|U_j^*|\leqslant u_b$、$|\Lambda_{U_j}^*|\leqslant Mu_b$（$j\in\mathbb{N}$）。因 $u_b\in L^\infty(0,T;L^1(\mathbb{R}))$，由控制收敛定理得
\begin{align}
\lim_{j\to\infty}U_j^*=\lim_{j\to\infty}\Lambda_{U_j}^*=0\quad\text{在 }S\text{ 中几乎处处且在 }L^1(S), \label{align:ch07-7.117}
\end{align}
从而在\eqref{align:ch07-7.116}中令 $j\to\infty$，对所有 $\tau\in(0,T)\setminus N$ 与 $\zeta\in C^1([0,T];C^1_c(\mathbb{R}))$ 得
\begin{align}
\langle\sigma(\cdot,\tau),\zeta(\cdot,\tau)\rangle
\leqslant\int_0^\tau\langle\sigma(\cdot,t),\partial_t\zeta(\cdot,t)\rangle\,\mathrm{d}t+\langle u_{0s},\zeta(\cdot,0)\rangle. \label{align:ch07-7.118}
\end{align}
在\eqref{align:ch07-7.118}中取 $\zeta(x,t)=\rho(x)$，得对任意非负 $\rho\in C^1_c(\mathbb{R})$ 有 $\langle\sigma(\cdot,\tau),\rho\rangle\leqslant\langle u_{0s},\rho\rangle$，故 $\sigma(\cdot,\tau)\ll u_{0s}$（回忆 $\sigma$ 非负）。于是对任意 $\tau\in(0,T)\setminus N$，$\sigma(\cdot,\tau)$ 关于 Lebesgue 测度奇异。另一方面，因 $u_b\in L^\infty(0,T;L^1(\mathbb{R}))$，测度 $E\mapsto\iint_E u_b\,\mathrm{d}x\,\mathrm{d}t$（$E\in\mathcal{B}(\mathbb{R}^2)\cap S$）关于 Lebesgue 测度绝对连续。由 Lebesgue 分解的唯一性（见\cref{theorem:ch01-Lebesgue分解}）得结论。
\end{proof}

% SUPPLEMENT: parent=proposition:ch07-7.5.5; kind=proof; location=命题 7.5.5 原证明中的三处省略; basis=原书对“将 (7.77) 与 (7.78) 相加”未写出 $U_j$ 与 $\Lambda_{U_j}$ 恰好是平均熵对（含常数项的消去）、“由标准逼近论证”未写出 (7.111) 的时间截断推导、由 (7.45) 到 (7.115) 未写出有界复合函数的弱$^*$极限识别、以及最后的 Lebesgue 分解只给结论
\begin{proof}
补出原证明中四处未展开的环节。

第一步：\eqref{align:ch07-7.110} 中的 $U_j,\Lambda_{U_j}$ 来自熵对与守恒律的平均。置 $E_j(y)=|y-j|$、$F_j(y)=\operatorname{sgn}(y-j)[\phi(y)-\phi(j)]$。直接计算给出
\begin{align*}
\frac{E_j(y)+y}{2}=(y-j)_+=U_j(y),\qquad
\frac{F_j(y)+\phi(y)}{2}=\frac{1+\operatorname{sgn}(y-j)}{2}[\phi(y)-\phi(j)]=\Lambda_{U_j}(y),
\end{align*}
其中用到 $\phi(y)=[\phi(y)-\phi(j)]+\phi(j)$ 与 $\operatorname{sgn}(y-j)+1=2\chi_{[j,\infty)}(y)$。对 $l=j$ 写出\eqref{align:ch07-7.77}，再与\eqref{align:ch07-7.78}相加并取平均，得
\begin{align*}
\iint_S\{U_j(u_{n_k})\partial_t\zeta+\Lambda_{U_j}(u_{n_k})\partial_x\zeta\}\,\mathrm{d}x\,\mathrm{d}t
+\int_{\mathbb{R}}U_j(u_{0n_k})\zeta(x,0)\,\mathrm{d}x
+\frac{j}{2}\iint_S\partial_t\zeta\,\mathrm{d}x\,\mathrm{d}t+\frac{\phi(j)}{2}\iint_S\partial_x\zeta\,\mathrm{d}x\,\mathrm{d}t\geqslant0.
\end{align*}
由 $\zeta(\cdot,T)=0$ 得 $\iint_S\partial_t\zeta\,\mathrm{d}x\,\mathrm{d}t=-\int_{\mathbb{R}}\zeta(x,0)\,\mathrm{d}x$；由 $\zeta(\cdot,t)\in C^1_c(\mathbb{R})$ 得 $\iint_S\partial_x\zeta\,\mathrm{d}x\,\mathrm{d}t=0$。故常数项 $\frac{j}{2}$ 的两项相消，$\frac{\phi(j)}{2}$ 的项为零，即得原证明开头的不等式。

第二步：\eqref{align:ch07-7.111} 的时间截断。设 $\zeta$ 如命题所述，$\zeta\geqslant0$，$\zeta(\cdot,T)=0$。对 $\delta\in(0,T-\tau)$ 取 $\psi_\delta\in C^1([0,T])$ 使 $\psi_\delta=1$ 于 $[0,\tau-\delta]$、$\psi_\delta=0$ 于 $[\tau,T]$、$\psi_\delta$ 线性于 $[\tau-\delta,\tau]$，并令 $\zeta_\delta\triangleq\zeta\psi_\delta$。把第一步的不等式用于 $\zeta_\delta$（它满足同样的假设：非负、$C^1$、在 $t=T$ 处为零）：
\begin{align*}
\iint_S\{U_j(u_{n_k})\partial_t\zeta_\delta+\Lambda_{U_j}(u_{n_k})\partial_x\zeta_\delta\}\,\mathrm{d}x\,\mathrm{d}t
\geqslant-\int_{\mathbb{R}}U_j(u_{0n_k})(x)\zeta_\delta(x,0)\,\mathrm{d}x.
\end{align*}
因 $\psi_\delta(0)=1$，右端与 $\delta$ 无关。又 $\partial_t\zeta_\delta=\psi_\delta\partial_t\zeta+\zeta\psi_\delta'$，$\partial_x\zeta_\delta=\psi_\delta\partial_x\zeta$。由 $\psi_\delta\to\chi_{(0,\tau)}$ 单调、$U_j(u_{n_k}),\Lambda_{U_j}(u_{n_k})\in L^1(S)$（见\eqref{align:ch07-7.71}与 $0\leqslant U_j(y)\leqslant y$），单调收敛定理给出前两项的极限为 $\iint_{S_\tau}$ 上的积分。对含 $\psi_\delta'$ 的项，
\begin{align*}
\iint_SU_j(u_{n_k})\zeta\psi_\delta'\,\mathrm{d}x\,\mathrm{d}t
=-\frac1\delta\int_{\tau-\delta}^{\tau}\!\!\int_{\mathbb{R}}U_j(u_{n_k})(x,t)\zeta(x,t)\,\mathrm{d}x\,\mathrm{d}t,
\end{align*}
而 $t\mapsto\int_{\mathbb{R}}U_j(u_{n_k})(x,t)\zeta(x,t)\,\mathrm{d}x\in L^1(0,T)$，故由 Lebesgue 微分定理，右端沿 $\delta\downarrow0$ 收敛于 $-\int_{\mathbb{R}}U_j(u_{n_k})(x,\tau)\zeta(x,\tau)\,\mathrm{d}x$ 对几乎所有 $\tau\in(0,T)$ 成立。合并即得\eqref{align:ch07-7.111}。

第三步：\eqref{align:ch07-7.115} 的识别。置 $h_j(y)\triangleq\min(y,j)$，则 $u_{0n_k}-U_j(u_{0n_k})=h_j(u_{0n_k})$。由\eqref{align:ch07-7.45}，$u_{0n_k}\to u_{0r}$ 于 $\mathbb{R}$ 中几乎处处且 $u_{0n_k}\stackrel{*}{\rightharpoonup}u_0$ 于 $\mathscr{R}_f(\mathbb{R})$；故与 $\{u_{0n_k}\}$ 相伴的 Young 测度在几乎每个 $x$ 处为 $\delta_{u_{0r}(x)}$。$h_j$ 有界连续，故由\cref{proposition:ch05-5.4.10}（取 $X$ 为 $\mathbb{R}$ 的紧子集、$h=h_j$，并用 $\zeta(\cdot,0)\in C_c(\mathbb{R})$ 的紧支撑）得
\begin{align*}
\lim_{k\to\infty}\int_{\mathbb{R}}h_j(u_{0n_k})(x)\zeta(x,0)\,\mathrm{d}x=\int_{\mathbb{R}}h_j(u_{0r})(x)\zeta(x,0)\,\mathrm{d}x.
\end{align*}
又由\eqref{align:ch07-7.45}，$\int_{\mathbb{R}}u_{0n_k}\zeta(x,0)\,\mathrm{d}x\to\langle u_0,\zeta(\cdot,0)\rangle=\int_{\mathbb{R}}u_{0r}\zeta(x,0)\,\mathrm{d}x+\langle u_{0s},\zeta(\cdot,0)\rangle$。两式相减，并用 $U_j(u_{0r})=u_{0r}-h_j(u_{0r})$，即得\eqref{align:ch07-7.115}。

第四步：\eqref{align:ch07-7.117} 与最后的 Lebesgue 分解。由\eqref{align:ch07-7.95}，
\begin{align*}
U_j^*(x,t)=\int_{[j,\infty)}(y-j)\,\mathrm{d}\nu_{(x,t)}(y)\leqslant\int_{[j,\infty)}y\,\mathrm{d}\nu_{(x,t)}(y),
\end{align*}
而后者对几乎所有 $(x,t)$ 随 $j\to\infty$ 单调下降至 $0$（因 $\int_{[0,\infty)}y\,\mathrm{d}\nu_{(x,t)}(y)=u_b(x,t)<\infty$）；由控制收敛定理（控制函数 $u_b\in L^1(S)$）得 $U_j^*\to0$ 于 $S$ 中几乎处处且在 $L^1(S)$。对 $\Lambda_{U_j}^*$ 作同样的论证，用 $|\Lambda_{U_j}(y)|\leqslant My\chi_{[j,\infty)}(y)$ 与控制函数 $Mu_b$。最后，由 $\sigma(\cdot,\tau)\ll u_{0s}$ 与 $u_{0s}\perp\lambda$（$\lambda$ 为 Lebesgue 测度）得 $\sigma(\cdot,\tau)\perp\lambda$ 对每个 $\tau\in(0,T)\setminus N$ 成立。取 $E\in\mathcal{B}(\mathbb{R}^2)\cap S$ 使 $\lambda_2(E)=0$ 且 $\sigma(\cdot,\tau)(E_\tau)=\|\sigma(\cdot,\tau)\|$ 对 $\tau\in(0,T)\setminus N$ 成立（$E_\tau$ 为 $t=\tau$ 处的截口），则 $\lambda_2(E)=0$（Fubini 定理），故 $\sigma\perp\lambda_2$；又 $u_b\,\mathrm{d}x\,\mathrm{d}t\ll\lambda_2$。于是 $u=u_b\,\mathrm{d}x\,\mathrm{d}t+\sigma$ 就是 $u$ 关于 $\lambda_2$ 的 Lebesgue 分解，由唯一性得\eqref{align:ch07-7.109a}--\eqref{align:ch07-7.109b}。
\end{proof}

% SUPPLEMENT: parent=proposition:ch07-7.5.5; kind=motivation; location=命题 7.5.5 后的注记; basis=命题 7.5.5 与 (7.109a)--(7.109b)、(7.116)、(7.118) + 为何 $U_j$ 取“截断到 $[j,\infty)$ 的一次函数”、$j\to\infty$ 极限与 $\tau$ 的零测集、$\sigma\ll u_{0s}$ 的来源与作用、以及该命题给出的分解为何恰是 Lebesgue 分解
\begin{remark}\phantomsection\label{remark:ch07-7.5.5动机}
\begin{enumerate}[(i)]
\item 引入动机。\cref{proposition:ch07-7.5.2}把极限测度写成 $u_b+\sigma$，但那只是“重心加集中”的分解，并未断言 $\sigma$ 奇异于 Lebesgue 测度（\cref{remark:ch07-7.5.2动机}(iv)）。\cref{proposition:ch07-7.5.5}用熵不等式补上这一信息：\eqref{align:ch07-7.109a}--\eqref{align:ch07-7.109b}说明 $u_b$ 恰是 $u$ 的绝对连续部分的密度、$\sigma$ 恰是奇异部分。这正是\cref{definition:ch07-熵解}中 $u$ 以 $u_r\,\mathrm{d}x\,\mathrm{d}t+u_s$ 形式出现的前提，也是\cref{theorem:ch07-7.5.1}(iii)中能把 $M_Uu_s$ 写成集中项的依据。
\item 截断函数的选择。取 $U_j(y)=(y-j)_+$ 而不是 $|y-j|$ 或 $y\wedge j$ 的原因有两点：其一，$U_j$ 的斜率 $\chi_{[j,\infty)}$ 在 $[0,\infty)$ 上单调不减，故 $U_j'\geqslant0$，从而 $\Lambda_{U_j}$ 可直接算出为 $[\phi(y)-\phi(j)]\chi_{[j,\infty)}(y)$；其二，$M_{U_j}=1$、$M_{\Lambda_{U_j}}=0$，使\eqref{align:ch07-7.112a}中只剩 $\langle\sigma,\partial_t\zeta\rangle$ 一项，而\eqref{align:ch07-7.112b}中没有 $\sigma$ 项——这正是\eqref{align:ch07-7.118}右端只含 $\sigma$ 与 $u_{0s}$ 的原因。若改用 $|y-j|$，则 $M_{U_j}=1$ 仍成立但 $\Lambda$ 的符号处理更繁；若改用 $y\wedge j$，则 $M_{U_j}=0$，得不到关于 $\sigma$ 的信息。
\item $j\to\infty$ 的极限与量词。\eqref{align:ch07-7.117}使 $U_j^*,\Lambda_{U_j}^*$ 消失，但\eqref{align:ch07-7.116}右端还有 $U_j(u_{0r})$，它随 $j\to\infty$ 单调下降至 $0$ 并被 $u_{0r}\in L^1(\mathbb{R})$ 控制，故也消失。剩下的\eqref{align:ch07-7.118}对每个固定的 $\tau\in(0,T)\setminus N$ 成立，其中零测集 $N$ 来自\eqref{align:ch07-7.107}的对角抽取，与 $j$、$\zeta$ 无关；这一点是必要的，因为最后要对所有 $\tau\notin N$ 断言 $\sigma(\cdot,\tau)$ 奇异。
\item $\sigma\ll u_{0s}$ 的来源与作用。在\eqref{align:ch07-7.118}中取与 $t$ 无关的 $\zeta=\rho$ 得 $\langle\sigma(\cdot,\tau),\rho\rangle\leqslant\langle u_{0s},\rho\rangle$ 对一切非负 $\rho\in C^1_c(\mathbb{R})$ 成立；由 $\sigma\geqslant0$ 与 $C^1_c(\mathbb{R})$ 在 $C_c(\mathbb{R})$ 中的稠密性，这等价于测度不等式 $\sigma(\cdot,\tau)\leqslant u_{0s}$，从而 $\sigma(\cdot,\tau)\ll u_{0s}$。由于 $u_{0s}$ 奇异于 Lebesgue 测度，得 $\sigma(\cdot,\tau)\perp\lambda$。这里 $\sigma\geqslant0$ 不可去掉：对带号测度，$\langle\sigma,\rho\rangle\leqslant\langle u_{0s},\rho\rangle$ 对非负 $\rho$ 成立只给出 $\sigma^+\leqslant u_{0s}$ 的一部分信息。
\item 边界。结论对几乎所有 $\tau$ 成立而不是对一切 $\tau$：$N$ 可能非空。此外证明只用到 $u_{0s}$ 的奇异性，故结论可陈述为“$\sigma$ 集中在一个 Lebesgue 零测集上”；但\eqref{align:ch07-7.109a}的等式只在几乎处处意义下成立，不能逐点断言。
\end{enumerate}
\end{remark}

\begin{proposition}\label{proposition:ch07-divcurl}\label{proposition:ch07-7.5.6}
设 $T>0$，$S\triangleq\mathbb{R}\times(0,T)$，(A0)--(A1) 成立（见\cref{subsection:ch07-假设}）。设 $\nu\in\mathcal{Y}^+(S;\mathbb{R})$ 由\cref{proposition:ch07-7.5.2}给出，$u\in L^\infty_{w^*}(0,T;\mathscr{R}_f^+(\mathbb{R}))$ 由\eqref{align:ch07-7.108}定义，$\{\nu_{(x,t)}\}_{(x,t)\in S}$ 为 $\nu$ 的离散分解。则
\begin{align}
\phi^*(x,t)=\phi(u_r(x,t))\quad\text{对几乎所有 }(x,t)\in S, \label{align:ch07-7.119}
\end{align}
其中 $\phi^*$ 由\eqref{align:ch07-7.4}定义。
\end{proposition}

\begin{proof}
设 $U,V\in C^2([0,\infty))\cap L^\infty(0,\infty)$ 满足\eqref{align:ch07-7.58}，并设 $\Lambda_U,\Lambda_V\in L^\infty(0,\infty)$（见\eqref{align:ch07-7.83}）。由\cref{lemma:ch07-7.4.8}，对任意 $n\in\mathbb{N}$，存在 $\sigma_n,\tau_n\in\mathscr{R}_f(S)$，使对所有 $\zeta\in C^1_c(S)$ 有
\begin{align}
\iint_S\{U(u_n)\partial_t\zeta+\Lambda_U(u_n)\partial_x\zeta\}\,\mathrm{d}x\,\mathrm{d}t=\langle\sigma_n,\zeta\rangle, \label{align:ch07-7.120a}
\end{align}
\begin{align}
\iint_S\{V(u_n)\partial_t\zeta+\Lambda_V(u_n)\partial_x\zeta\}\,\mathrm{d}x\,\mathrm{d}t=\langle\tau_n,\zeta\rangle, \label{align:ch07-7.120b}
\end{align}
其中 $u_n$ 为\cref{proposition:ch07-7.4.7}给出的 (Pn) 的熵解。

设 $A\Subset S$ 为有界开集，定义 $Y_n,Z_n:A\to\mathbb{R}^2$ 为
\begin{align*}
Y_n\triangleq (\Lambda_U(u_n),U(u_n)),\qquad Z_n\triangleq (V(u_n),-\Lambda_V(u_n)).
\end{align*}
由\eqref{align:ch07-7.120a}得
\begin{align}
\operatorname{div}Y_n=-\sigma_n,\qquad \operatorname{curl}Z_n=-\tau_n\quad\text{在 }\mathcal{D}^*(A). \label{align:ch07-7.121}
\end{align}
因 $U,\Lambda_U,V,\Lambda_V$ 在 $(0,\infty)$ 上有界，序列 $\{U(u_n)\}$、$\{\Lambda_U(u_n)\}$、$\{V(u_n)\}$、$\{\Lambda_V(u_n)\}$ 在 $L^\infty(A)$ 中有界，从而在 $L^1(A)$ 中有界且一致可积。由\cref{proposition:ch05-5.4.10}，在 $L^1(A)$ 中
\begin{align*}
U(u_n)\rightharpoonup U^*\triangleq \int_{[0,\infty)}U(y)\,\mathrm{d}\nu(y),\quad
\Lambda_U(u_n)\rightharpoonup\Lambda_U^*\triangleq \int_{[0,\infty)}\Lambda_U(y)\,\mathrm{d}\nu(y),
\end{align*}
\begin{align*}
V(u_n)\rightharpoonup V^*\triangleq \int_{[0,\infty)}V(y)\,\mathrm{d}\nu(y),\quad
\Lambda_V(u_n)\rightharpoonup\Lambda_V^*\triangleq \int_{[0,\infty)}\Lambda_V(y)\,\mathrm{d}\nu(y).
\end{align*}
由同样的论证，$\{U(u_n)\}$、$\{\Lambda_U(u_n)\}$、$\{V(u_n)\}$、$\{\Lambda_V(u_n)\}$ 在 $L^\infty(A)$ 中有界，从而在 $L^2(A)$ 中弱收敛。于是当 $n\to\infty$ 时
\begin{align}
Y_n\rightharpoonup Y^*\triangleq (\Lambda_U^*,U^*),\quad
Z_n\rightharpoonup Z^*\triangleq (V^*,-\Lambda_V^*)\quad\text{在 }L^2(A;\mathbb{R}^2). \label{align:ch07-7.122}
\end{align}
由同样的论证可得
\begin{align}
Y_n\cdot Z_n\triangleq \Lambda_U(u_n)V(u_n)-\Lambda_V(u_n)U(u_n)
\rightharpoonup\int_{[0,\infty)}[\Lambda_U(y)V(y)-\Lambda_V(y)U(y)]\,\mathrm{d}\nu(y)
\quad\text{在 }L^2(A). \label{align:ch07-7.123}
\end{align}
由\eqref{align:ch07-7.85}（应用于 $\sigma_n$ 与 $\tau_n$）与\eqref{align:ch07-7.121}，$\{\operatorname{div}Y_n\}$ 与 $\{\operatorname{curl}Z_n\}$ 在 $H^{-1}(A)$ 中预紧。于是由\eqref{align:ch07-7.122}与 div--curl 引理（Murat--Tartar）得
\begin{align}
Y_n\cdot Z_n\to Y^*\cdot Z^*=\Lambda_U^*V^*-\Lambda_V^*U^*\quad\text{在 }\mathcal{D}^*(A) \label{align:ch07-7.124}
\end{align}
（例如见文献~\cite{ref46} 推论 1.3.1 与\cref{theorem:ch05-离散分解}。由\eqref{align:ch07-7.123}--\eqref{align:ch07-7.124}得
\begin{align}
\int_{[0,\infty)}[\Lambda_U(y)-\Lambda_U^*]V(y)\,\mathrm{d}\nu(y)
=\int_{[0,\infty)}[U(y)-U^*]\Lambda_V(y)\,\mathrm{d}\nu(y)
\quad\text{在 }A\text{ 中几乎处处}. \label{align:ch07-7.125}
\end{align}

设 $U$ 如上，$U'>0$ 在 $(0,\infty)$。由标准逼近论证可取 $V(y)=|U(y)-U^*|$，适当选取 $\vartheta_V$ 得
\begin{align*}
\Lambda_V(y)=\int_0^y V'(z)\phi'(z)\,\mathrm{d}z+\vartheta_V
=\operatorname{sgn}(U(y)-U^*)[\Lambda_U(y)-\Lambda_U(U^{-1}(U^*))],
\end{align*}
故由\eqref{align:ch07-7.125}得
\begin{align}
[\Lambda_U^*-\Lambda_U(U^{-1}(U^*))]\int_{[0,\infty)}|U(y)-U^*|\,\mathrm{d}\nu(y)=0
\quad\text{在 }A\text{ 中几乎处处}. \label{align:ch07-7.126}
\end{align}
设 $U_p\in C^2([0,\infty))\cap L^\infty(0,\infty)$、$U_p(0)=0$（$p\in\mathbb{N}$），满足\eqref{align:ch07-7.58}及
\begin{align}
0<U_p'\leqslant U_{p+1}'\leqslant1,\quad \lim_{p\to\infty}U_p'(y)=1\ \text{在 }[0,\infty), \label{align:ch07-7.127}
\end{align}
从而 $\lim_{p\to\infty}U_p(y)=y$ 对所有 $y\in[0,\infty)$ 成立。由单调收敛定理、\eqref{align:ch07-7.95}与\eqref{align:ch07-7.109a}得
\begin{align}
\lim_{p\to\infty}U_p^*=\lim_{p\to\infty}\int_{[0,\infty)}U_p(y)\,\mathrm{d}\nu(y)=u_r\quad\text{在 }A\text{ 中几乎处处}. \label{align:ch07-7.128}
\end{align}
下面证明
\begin{align}
\lim_{p\to\infty}[\Lambda_{U_p}^*-\Lambda_{U_p}(U_p^{-1}(U_p^*))]=\phi^*-\phi(u_r)\quad\text{在 }A\text{ 中几乎处处}. \label{align:ch07-7.129}
\end{align}
利用\eqref{align:ch07-7.128}--\eqref{align:ch07-7.129}即可完成证明。事实上，由\eqref{align:ch07-7.128}与控制收敛定理，对几乎所有 $(x,t)\in A$ 有
\begin{align}
\lim_{p\to\infty}\int_{[0,\infty)}|U_p(y)-U_p^*(x,t)|\,\mathrm{d}\nu_{(x,t)}(y)
=\int_{[0,\infty)}|y-u_r(x,t)|\,\mathrm{d}\nu_{(x,t)}(y). \label{align:ch07-7.130}
\end{align}
在\eqref{align:ch07-7.126}（取 $U=U_p$）中令 $p\to\infty$，由\eqref{align:ch07-7.129}--\eqref{align:ch07-7.130}，对几乎所有 $(x,t)\in A$ 得
\begin{align*}
[\phi^*(x,t)-\phi(u_r(x,t))]\int_{[0,\infty)}|y-u_r(x,t)|\,\mathrm{d}\nu_{(x,t)}(y)=0.
\end{align*}
由 $A$ 的任意性，上式蕴含\eqref{align:ch07-7.119}。

余下需证\eqref{align:ch07-7.129}。直接计算给出
\begin{align}
\Lambda_{U_p}^*-\Lambda_{U_p}(U_p^{-1}(U_p^*))
=\int_{[0,\infty)}\left(\int_0^y U_p'(z)\phi'(z)\,\mathrm{d}z\right)\mathrm{d}\nu(y)
-\int_0^{U_p^{-1}(U_p^*)}U_p'(z)\phi'(z)\,\mathrm{d}z. \label{align:ch07-7.131}
\end{align}
由\eqref{align:ch07-7.1}与\eqref{align:ch07-7.127}，$U_p'(y)\to1$ 且 $|U_p'(y)\phi'(y)|\leqslant M$ 对所有 $y\geqslant0$ 成立，故由控制收敛定理
\begin{align}
\lim_{p\to\infty}\Lambda_{U_p}^*(x,t)=\phi^*(x,t)\quad\text{对几乎所有 }(x,t)\in A. \label{align:ch07-7.132}
\end{align}
另一方面，
\begin{align}
\int_0^{U_p^{-1}(U_p^*(x,t))}U_p'(z)\phi'(z)\,\mathrm{d}z-\phi(u_r)(x,t)
=\int_0^{u_r(x,t)}[U_p'(z)-1]\phi'(z)\,\mathrm{d}z
+\int_{u_r(x,t)}^{U_p^{-1}(U_p^*(x,t))}U_p'(z)\phi'(z)\,\mathrm{d}z. \label{align:ch07-7.133}
\end{align}
如\eqref{align:ch07-7.132}那样论证，得
\begin{align}
\lim_{p\to\infty}\int_0^{u_r(x,t)}[U_p'(z)-1]\phi'(z)\,\mathrm{d}z=0
\quad\text{对几乎所有 }(x,t)\in A. \label{align:ch07-7.134}
\end{align}
关于\eqref{align:ch07-7.133}右端第二项，由\eqref{align:ch07-7.127}，对充分大的 $p$ 有
\begin{align*}
\left|\int_{u_r(x,t)}^{U_p^{-1}(U_p^*(x,t))}U_p'(z)\phi'(z)\,\mathrm{d}z\right|
\leqslant M|u_r(x,t)-U_p^{-1}(U_p^*(x,t))|,
\end{align*}
故由\eqref{align:ch07-7.128}令 $p\to\infty$ 得
\begin{align}
\lim_{p\to\infty}\int_{u_r(x,t)}^{U_p^{-1}(U_p^*(x,t))}U_p'(z)\phi'(z)\,\mathrm{d}z=0
\quad\text{对几乎所有 }(x,t)\in A. \label{align:ch07-7.135}
\end{align}
由\eqref{align:ch07-7.133}--\eqref{align:ch07-7.135}得
\begin{align}
\int_0^{U_p^{-1}(U_p^*(x,t))}U_p'(z)\phi'(z)\,\mathrm{d}z\to\phi(u_r)(x,t)
\quad\text{对几乎所有 }(x,t)\in A. \label{align:ch07-7.136}
\end{align}
由\eqref{align:ch07-7.131}、\eqref{align:ch07-7.132}与\eqref{align:ch07-7.136}得\eqref{align:ch07-7.129}。证毕。
\end{proof}

% SUPPLEMENT: parent=proposition:ch07-7.5.6; kind=proof; location=命题 7.5.6 原证明中“由标准逼近论证可取 $V(y)=|U(y)-U^*|$”与 (7.125)->(7.126) 的推导、以及 (7.135) 的估计; basis=原书以“standard approximation arguments”“by a proper choice of $\vartheta_V$”带过两处、p.354 处 (7.135) 前的估计在抽取文本中残缺；此处补出 $V$ 的选取（含 $c=U^*$ 冻结与磨光）、$\vartheta_V$ 的显式计算、(7.126) 的代数推导与 (7.135) 的反函数估计
\begin{proof}
补出原证明中与 $V$ 的选取有关的三处环节。

第一步：$V(y)=|U(y)-U^*|$ 的合法性。(7.125) 的证明只对固定的 $V\in C^2([0,\infty))\cap L^\infty(0,\infty)$ 进行，而 $|U(\cdot)-U^*|$ 通过 $U^*(x,t)$ 依赖 $(x,t)$。取 $V_c(y)\triangleq|U(y)-c|$（$c\in\mathbb{R}$ 固定）。因 $U'>0$，$V_c$ 在 $U^{-1}(c)$ 处有折点，故先用
\begin{align*}
V_{c,\delta}(y)\triangleq\sqrt{(U(y)-c)^2+\delta^2}-\delta\quad(\delta>0)
\end{align*}
代替。为使\cref{lemma:ch07-7.4.8}的前提\eqref{align:ch07-7.58}也得到满足，把 $V_{c,\delta}$ 在 $[0,R]$ 之外改为线性延拓并磨光，所得函数在 $[0,R]$ 上 $V''$ 有界、在 $[R,\infty)$ 上 $V''=0$，故\eqref{align:ch07-7.58}成立；再令 $R\to\infty$ 与 $\delta\downarrow0$。两次取极限都由下列与 $c,\delta,R$ 无关的界控制：由 $U'>0$ 得 $\int_0^y|U'|=U(y)-U(0)$，故 $|\Lambda_{V_{c,\delta}}(y)|\leqslant M\int_0^y|U'|=M[U(y)-U(0)]\leqslant2M\|U\|_{L^\infty(0,\infty)}$，又 $|V_{c,\delta}|\leqslant|U-c|\leqslant\|U\|_{L^\infty}+|c|$。因 $\nu_{(x,t)}$ 是概率测度，上述常数即为 $y$-积分意义下的优函数，故控制收敛定理给出\eqref{align:ch07-7.125}对 $V=V_c$ 成立。该等式对每个固定的 $c\in\mathbb{R}$ 成立，且两端关于 $c$ 连续，故在 $c=U^*(x,t)$（$(x,t)$ 的可测函数）处取值仍成立，即得\eqref{align:ch07-7.125}中取 $V(y)=|U(y)-U^*|$ 的情形。

第二步：$\vartheta_V$ 的显式计算。设 $U'>0$、$\alpha\triangleq U^{-1}(U^*)$。由\eqref{align:ch07-7.83}，$\Lambda_U(y)=\int_0^yU'(z)\phi'(z)\,\mathrm{d}z+\vartheta_U$，且 $V'(z)=\operatorname{sgn}(U(z)-U^*)U'(z)$。故
\begin{align*}
\int_0^yV'(z)\phi'(z)\,\mathrm{d}z=\int_0^y\operatorname{sgn}(U(z)-U^*)U'(z)\phi'(z)\,\mathrm{d}z
=\operatorname{sgn}(U(y)-U^*)\int_\alpha^yU'(z)\phi'(z)\,\mathrm{d}z
\end{align*}
（当 $y>\alpha$ 时符号恒为 $+1$，当 $y<\alpha$ 时恒为 $-1$）。而右端等于 $\operatorname{sgn}(U(y)-U^*)[\Lambda_U(y)-\Lambda_U(\alpha)]$。再取
\begin{align*}
\vartheta_V\triangleq-\operatorname{sgn}(U(0)-U^*)\big[\Lambda_U(0)-\Lambda_U(\alpha)\big]
\end{align*}
即可使 $\Lambda_V(y)=\int_0^yV'\phi'+\vartheta_V$ 与 $\operatorname{sgn}(U(y)-U^*)[\Lambda_U(y)-\Lambda_U(\alpha)]$ 在所有 $y$ 处相等（两者只差一个常数，而后者本身连续）。此即“适当选取 $\vartheta_V$”的确切含义。

第三步：由\eqref{align:ch07-7.125}到\eqref{align:ch07-7.126}。以 $V(y)=|U(y)-U^*|$ 代入\eqref{align:ch07-7.125}，左端为 $\int[\Lambda_U(y)-\Lambda_U^*]|U(y)-U^*|\,\mathrm{d}\nu(y)$；右端由第二步为 $\int[U(y)-U^*]\operatorname{sgn}(U(y)-U^*)[\Lambda_U(y)-\Lambda_U(\alpha)]\,\mathrm{d}\nu(y)=\int|U(y)-U^*|[\Lambda_U(y)-\Lambda_U(\alpha)]\,\mathrm{d}\nu(y)$。两式相减得
\begin{align*}
\int_{[0,\infty)}|U(y)-U^*|\big\{[\Lambda_U(y)-\Lambda_U^*]-[\Lambda_U(y)-\Lambda_U(\alpha)]\big\}\,\mathrm{d}\nu(y)=0,
\end{align*}
即 $[\Lambda_U(\alpha)-\Lambda_U^*]\int|U(y)-U^*|\,\mathrm{d}\nu(y)=0$，这正是\eqref{align:ch07-7.126}。

第四步：\eqref{align:ch07-7.135} 的反函数估计。由\eqref{align:ch07-7.127}，$U_p$ 严格递增且 $U_p'\geqslant U_1'>0$ 在 $(0,\infty)$ 上。设 $u_r=u_r(x,t)>0$，取 $\delta\in(0,u_r/2)$。对充分大的 $p$，$U_p^*(x,t)\in(u_r-\delta,u_r+\delta)$（见\eqref{align:ch07-7.128}）且 $U_p^{-1}(u_r)$ 与 $u_r$ 之差不超过 $\delta$。由中值定理与 $U_p^{-1}$ 的 Lipschitz 估计，
\begin{align*}
|U_p^{-1}(U_p^*)-u_r|\leqslant|U_p^{-1}(u_r)-u_r|+\frac{1}{\inf_{z\in(u_r-\delta,u_r+\delta)}U_1'(z)}|u_r-U_p^*|,
\end{align*}
故
\begin{align*}
\left|\int_{u_r}^{U_p^{-1}(U_p^*)}U_p'(z)\phi'(z)\,\mathrm{d}z\right|
&\leqslant M|u_r-U_p^{-1}(U_p^*)|
\leqslant M\Big\{|u_r-U_p^{-1}(u_r)|+\Big(\inf_{z\in(u_r-\delta,u_r+\delta)}U_1'(z)\Big)^{-1}|u_r-U_p^*|\Big\}.
\end{align*}
因 $U_p(y)\to y$ 且 $U_p,U_p^{-1}$ 在紧集上一致收敛于恒等映射，得 $U_p^{-1}(u_r)\to u_r$；又由\eqref{align:ch07-7.128}得 $U_p^*\to u_r$。于是右端两项均趋于 $0$，即\eqref{align:ch07-7.135}。
\end{proof}

% SUPPLEMENT: parent=proposition:ch07-7.5.6; kind=motivation; location=命题 7.5.6 后的注记; basis=命题 7.5.6 与 (7.119)、(7.125)--(7.129) + div-curl 补偿紧性的动机、为何需要两对熵对、$U'>0$ 的作用、$\{U_p\}$ 的选取、以及该命题不需要假设 (A2) 而 (A2) 只在命题 7.5.7 中进入的边界
\begin{remark}\phantomsection\label{remark:ch07-7.5.6动机}
\begin{enumerate}[(i)]
\item 引入动机。\cref{theorem:ch07-7.5.1}(ii)只给出 $\phi(u_r)=\int\phi\,\mathrm{d}\nu$，这不足以断言 $\phi(u_r)$ 就是 $\phi$ 在振荡分布上的“正确取值”：若 $\nu_{(x,t)}$ 非退化，则 $\int\phi\,\mathrm{d}\nu\neq\phi(\int y\,\mathrm{d}\nu)$（除非 $\phi$ 在 $\operatorname{supp}\nu_{(x,t)}$ 上仿射）。\cref{proposition:ch07-7.5.6}用 div-curl 引理证明\eqref{align:ch07-7.119}，即 $\phi$ 与 Young 测度“交换”：$\phi$ 在 $\nu_{(x,t)}$ 上的平均值等于它在重心处的值。这一步是\cref{theorem:ch07-7.3.2}的关键，因为它把 Young 测度熵解的不等式\eqref{align:ch07-7.5}中的 $\phi^*$ 换成了 $\phi(u_r)$，即换成了\cref{definition:ch07-Qtau熵解}中出现的量。
\item 为何需要两对熵对。\eqref{align:ch07-7.119}的证明是 Tartar 的补偿紧性论证：取两个熵对 $(U,\Lambda_U)$ 与 $(V,\Lambda_V)$，由 (Pn) 的熵等式把它们的通量写成散度形式\eqref{align:ch07-7.121}；div-curl 引理说明两个弱收敛向量场的点积在 $\mathcal{D}^*$ 中收敛（\eqref{align:ch07-7.124}），而\eqref{align:ch07-7.123}又给出该点积的另一种表达式，两者相等即得\eqref{align:ch07-7.125}。这里的核心是：单个熵对只给出一条守恒律，其散度只在一个方向上受控；必须有两个熵对，才能同时控制 $\operatorname{div}$ 与 $\operatorname{curl}$。
\item $U'>0$ 的作用与 $\{U_p\}$ 的选取。\eqref{align:ch07-7.126}中出现 $U^{-1}(U^*)$，故 $U$ 必须严格递增，这是假设 $U'>0$ 的唯一用途；它保证 $\Lambda_U(U^{-1}(U^*))$ 有定义。为了从\eqref{align:ch07-7.126}过渡到\eqref{align:ch07-7.119}，需要 $U$ 同时逼近恒等映射以使 $U^*\to u_r$（\eqref{align:ch07-7.128}）且 $\Lambda_U(U^{-1}(U^*))\to\phi(u_r)$（\eqref{align:ch07-7.136}）；这正是\eqref{align:ch07-7.127}中 $0<U_p'\leqslant U_{p+1}'\leqslant1$ 与 $U_p'\to1$ 的作用：单调性使 $U_p^*\to u_r$ 可用单调收敛定理，$U_p'\leqslant1$ 使 $|U_p^{-1}(a)-U_p^{-1}(b)|\geqslant|a-b|$ 方向的反向估计成立（反函数 Lipschitz 常数的界）。
\item 假设的作用。(A0)--(A1) 只通过\eqref{align:ch07-7.71}与\eqref{align:ch07-7.1}进入：前者给出 $u_n$ 的 $L^\infty(S)$ 与 $L^\infty(0,T;L^1)$ 界，使\cref{proposition:ch05-5.4.10}可用；后者给出 $|\phi'|\leqslant M$，使\eqref{align:ch07-7.132}、\eqref{align:ch07-7.134}、\eqref{align:ch07-7.135}中的控制函数可取 $M$。\eqref{align:ch07-7.58}与 $\Lambda_U,\Lambda_V\in L^\infty(0,\infty)$ 是\cref{lemma:ch07-7.4.8}的前提：$\{U(u_n)\}$ 等在 $L^\infty(A)$ 中有界，从而在 $L^2(A)$ 中弱收敛。
\item 边界。\cref{proposition:ch07-7.5.6}不需要假设 (A2)：它对任意 $\phi$ 给出\eqref{align:ch07-7.119}。(A2)（$\phi'$ 的严格单调性）只在\cref{proposition:ch07-7.5.7}中进入，用于证明 $\nu_{(x,t)}$ 是 Dirac 测度。反过来，若 $\phi'$ 在某区间上为常数，则\eqref{align:ch07-7.119}仍成立而 $\nu_{(x,t)}$ 可以非退化，见\cref{remark:ch07-7.5.3}。
\end{enumerate}
\end{remark}

\begin{proof}[\cref{theorem:ch07-7.5.1}的证明]
设 $u\in L^\infty_{w^*}(0,T;\mathscr{R}_f(\mathbb{R}))$ 由\eqref{align:ch07-7.108}定义。则由\cref{proposition:ch07-7.5.2}--\cref{proposition:ch07-7.5.6}得结论。
\end{proof}

% SUPPLEMENT: parent=theorem:ch07-7.5.1; kind=proof; location=定理 7.5.1 原证明后的补证; basis=原书仅写“由命题 7.5.2--7.5.6 得结论”，未指明哪一命题给出哪一项；此处逐项对应 (i)(ii)(iii) 与 (7.88)、(7.89)--(7.90)、(7.92) 的来源
\begin{proof}
补出三项陈述各自的依据。

第 (i) 项。由\cref{proposition:ch07-7.5.2}的\eqref{align:ch07-7.94}与\cref{proposition:ch07-7.5.5}的\eqref{align:ch07-7.108}，$u\triangleq u_b+\sigma\in L^\infty_{w^*}(0,T;\mathscr{R}_f^+(\mathbb{R}))$ 且 $u_{n_k}\stackrel{*}{\rightharpoonup}u$ 于 $L^\infty_{w^*}(0,T;\mathscr{R}_f(\mathbb{R}))$，即\eqref{align:ch07-7.88}。

第 (ii) 项。\cref{proposition:ch07-7.5.2}给出 Young 测度 $\nu\in\mathcal{Y}^+(S;\mathbb{R})$ 与离散分解 $\{\nu_{(x,t)}\}$；窄收敛性由\cref{theorem:ch05-Young弱星紧}与\cref{definition:ch05-窄收敛}（子列按弱$^*$收敛选取，且 $\|\nu_{n_k}\|=\|u_{n_k}\|_{L^1(S)}$ 有界、$\nu_{n_k}\geqslant0$，故弱$^*$极限是 $\mathcal{Y}^+$ 中的元素即窄极限）。\eqref{align:ch07-7.89}由\eqref{align:ch07-7.95}与\eqref{align:ch07-7.109a}直接得到；\eqref{align:ch07-7.90}由\eqref{align:ch07-7.119}与\eqref{align:ch07-7.4}得到。

第 (iii) 项。由\cref{proposition:ch07-7.5.3}的\eqref{align:ch07-7.98}，把其中 $\sigma$ 按\eqref{align:ch07-7.109b}换成 $u_s$ 即得\eqref{align:ch07-7.92}；$\nu$ 与 $u_s$ 均由同一子列给出，故 $U^*$ 的定义\eqref{align:ch07-7.93}与 (ii) 中的 $\nu$ 一致。
\end{proof}

现可证明\cref{theorem:ch07-7.3.1}。

\begin{proof}[\cref{theorem:ch07-7.3.1}的证明]
证明由\eqref{align:ch07-7.108}定义的 $u\in L^\infty_{w^*}(0,T;\mathscr{R}_f(\mathbb{R}))$ 满足等式
\begin{align}
\iint_S[u_r\partial_t\zeta+\phi(u_r)\partial_x\zeta]\,\mathrm{d}x\,\mathrm{d}t
+\int_0^T\langle u_s(\cdot,t),\partial_t\zeta(\cdot,t)\rangle\,\mathrm{d}t
=-\langle u_0,\zeta(\cdot,0)\rangle \label{align:ch07-7.137}
\end{align}
对所有 $\zeta\in C^1([0,T];C^1_c(\mathbb{R}))$、$\zeta(\cdot,T)=0$ 在 $\mathbb{R}$ 上成立，从而 $u$ 是 (P) 的解（见\cref{definition:ch07-熵解}）。设 $\{u_{n_k}\}\subseteq\{u_n\}$ 为\cref{theorem:ch07-7.5.1}给出的子列。对 $U(y)=y$ 与 $U(y)=\phi(y)$ 应用\cref{theorem:ch07-7.5.1}(iii)，得
\begin{align}
\lim_{k\to\infty}\iint_S u_{n_k}\partial_t\zeta\,\mathrm{d}x\,\mathrm{d}t
=\iint_S u_r\partial_t\zeta\,\mathrm{d}x\,\mathrm{d}t+\int_0^T\langle u_s(\cdot,t),\partial_t\zeta(\cdot,t)\rangle\,\mathrm{d}t, \label{align:ch07-7.138a}
\end{align}
\begin{align}
\lim_{k\to\infty}\iint_S\phi(u_{n_k})\partial_x\zeta\,\mathrm{d}x\,\mathrm{d}t
=\iint_S\phi(u_r)\partial_x\zeta\,\mathrm{d}x\,\mathrm{d}t \label{align:ch07-7.138b}
\end{align}
（见\eqref{align:ch07-7.89}、\eqref{align:ch07-7.4}与\eqref{align:ch07-7.119}）。由\eqref{align:ch07-7.45}中第一个收敛与\eqref{align:ch07-7.138a}，在\eqref{align:ch07-7.78}（取 $n=n_k$）中令 $k\to\infty$ 得\eqref{align:ch07-7.137}。不等式\eqref{align:ch07-7.5}由同样的取极限论证得到：在\eqref{align:ch07-7.77}（取 $n=n_k$）中令 $k\to\infty$，对 $U(y)=E_l(y)$ 与 $U(y)=F_l(y)$ 应用\cref{theorem:ch07-7.5.1}(iii)，并如\cref{proposition:ch07-7.5.3}证明中那样论证 $E_l(u_{0n_k})\stackrel{*}{\rightharpoonup}E_l(u_{0r})+u_{0s}$ 在 $\mathscr{R}_f(\mathbb{R})$ 中。结论得证。
\end{proof}

% SUPPLEMENT: parent=theorem:ch07-7.3.1; kind=proof; location=定理 7.3.1 原证明后的补证; basis=原证明以“对 $U(y)=y$ 与 $U(y)=\phi(y)$ 应用定理 7.5.1(iii)”带过对 (7.91) 的核验与 $M_U$ 的取值，又以“类似地”带过熵不等式 (7.5) 中 $E_l$、$F_l$ 对应的 $M_U$ 与边界项；此处补出两处核验与 (7.5) 的取极限
\begin{proof}
补出两处取极限的全部核验。

第一步：\eqref{align:ch07-7.138a}。取 $U(y)=y$。则 $U\in C([0,\infty))$ 且
\begin{align*}
M_U=\lim_{y\to\infty}\frac{y}{y}=1,\qquad
U^*(x,t)=\int_{[0,\infty)}y\,\mathrm{d}\nu_{(x,t)}(y)=u_r(x,t)
\end{align*}
（见\eqref{align:ch07-7.89}）。故\eqref{align:ch07-7.92}即\eqref{align:ch07-7.138a}；其中 $\partial_t\zeta\in C([0,T];C_c(\mathbb{R}))$，因为 $\zeta\in C^1([0,T];C^1_c(\mathbb{R}))$。

第二步：\eqref{align:ch07-7.138b}。取 $U(y)=\phi(y)$。由 (A1)，$M_\phi=\lim_{y\to\infty}\phi(y)/y=0$；又 $U^*=\int\phi\,\mathrm{d}\nu=\phi^*$（见\eqref{align:ch07-7.4}），而由\eqref{align:ch07-7.119}得 $\phi^*=\phi(u_r)$。故\eqref{align:ch07-7.92}中 $M_U$ 项消失，即\eqref{align:ch07-7.138b}。

第三步：\eqref{align:ch07-7.137}。对 $\zeta\in C^1([0,T];C^1_c(\mathbb{R}))$、$\zeta(\cdot,T)=0$，在\eqref{align:ch07-7.78}（取 $n=n_k$）中有
\begin{align*}
\iint_S[u_{n_k}\partial_t\zeta+\phi(u_{n_k})\partial_x\zeta]\,\mathrm{d}x\,\mathrm{d}t+\int_{\mathbb{R}}u_{0n_k}(x)\zeta(x,0)\,\mathrm{d}x=0.
\end{align*}
由\eqref{align:ch07-7.138a}--\eqref{align:ch07-7.138b}，前两项的极限为\eqref{align:ch07-7.137}的左端；由\eqref{align:ch07-7.45}中第一个收敛 $u_{0n_k}\stackrel{*}{\rightharpoonup}u_0$ 于 $\mathscr{R}_f(\mathbb{R})$ 与 $\zeta(\cdot,0)\in C^1_c(\mathbb{R})$，边界项收敛于 $\langle u_0,\zeta(\cdot,0)\rangle$。故\eqref{align:ch07-7.137}成立。

第四步：熵不等式\eqref{align:ch07-7.5}。设 $l\in[0,\infty)$、$\zeta\geqslant0$。在\eqref{align:ch07-7.77}（取 $n=n_k$）中，
\begin{align*}
\iint_S[E_l(u_{n_k})\partial_t\zeta+F_l(u_{n_k})\partial_x\zeta]\,\mathrm{d}x\,\mathrm{d}t+\int_{\mathbb{R}}E_l(u_{0n_k})\zeta(x,0)\,\mathrm{d}x\geqslant0.
\end{align*}
对 $U=E_l$：$M_{E_l}=\lim_{y\to\infty}|y-l|/y=1$（$l\geqslant0$），$E_l^*=\int|y-l|\,\mathrm{d}\nu$，故由\eqref{align:ch07-7.92}得
\begin{align*}
\lim_{k\to\infty}\iint_SE_l(u_{n_k})\partial_t\zeta\,\mathrm{d}x\,\mathrm{d}t
=\iint_SE_l^*\partial_t\zeta\,\mathrm{d}x\,\mathrm{d}t+\int_0^T\langle u_s(\cdot,t),\partial_t\zeta(\cdot,t)\rangle\,\mathrm{d}t.
\end{align*}
对 $U=F_l$：由\eqref{align:ch07-7.1}得 $|F_l(y)|\leqslant|\phi(y)-\phi(l)|\leqslant M(y+l)$，故 $M_{F_l}=\lim_{y\to\infty}F_l(y)/y=0$，且 $F_l^*=\int F_l\,\mathrm{d}\nu$，于是
\begin{align*}
\lim_{k\to\infty}\iint_SF_l(u_{n_k})\partial_x\zeta\,\mathrm{d}x\,\mathrm{d}t=\iint_SF_l^*\partial_x\zeta\,\mathrm{d}x\,\mathrm{d}t.
\end{align*}
最后处理边界项。因 $u_{0n_k}\geqslant0$，把 $E_l$ 延拓为 $\mathbb{R}$ 上的连续函数并取 $y\to-\infty$ 时的斜率为 $0$，则由\cref{proposition:ch05-5.4.13}（取 $X$ 为 $\mathbb{R}$ 的紧子集）与\eqref{align:ch07-7.45}得 $E_l(u_{0n_k})\stackrel{*}{\rightharpoonup}E_l(u_{0r})+1\cdot u_{0s}$ 于 $\mathscr{R}_f(\mathbb{R})$，即
\begin{align*}
\lim_{k\to\infty}\int_{\mathbb{R}}E_l(u_{0n_k})\zeta(x,0)\,\mathrm{d}x=\int_{\mathbb{R}}E_l(u_{0r})\zeta(x,0)\,\mathrm{d}x+\langle u_{0s},\zeta(\cdot,0)\rangle.
\end{align*}
合并三式即得\eqref{align:ch07-7.5}。由于 $l$ 与 $\zeta\geqslant0$ 任意，$u$ 是 (P) 的 Young 测度熵解；结合第三步得\cref{theorem:ch07-7.3.1}的结论。
\end{proof}

为证\cref{theorem:ch07-7.3.2}，需要下列结果，它刻画 Young 测度 $\nu$ 的离散分解。

\begin{proposition}\label{proposition:ch07-7.5.7}
设 $T>0$，$S\triangleq\mathbb{R}\times(0,T)$，(A0)--(A2) 成立（见\cref{subsection:ch07-假设}；特别地 $\phi\in C^1([0,\infty))$ 且对每个 $y_1>0$ 存在 $a,b\geqslant0$、$a+b>0$ 使 $\phi'$ 在 $[y_1-a,y_1+b]$ 上严格单调）。设 $u\in L^\infty_{w^*}(0,T;\mathscr{R}_f^+(\mathbb{R}))$ 与 $\nu\in\mathcal{Y}^+(S;\mathbb{R})$ 由\cref{theorem:ch07-7.5.1}给出，$\{\nu_{(x,t)}\}_{(x,t)\in S}$ 为 $\nu$ 的离散分解。则对几乎所有 $(x,t)\in S$ 有
\begin{align}
\nu_{(x,t)}=\delta_{u_r(x,t)}, \label{align:ch07-7.139}
\end{align}
其中 $u_r$ 为 $u$ 关于 Lebesgue 测度的绝对连续部分的密度，$\delta_{y}$ 为集中在 $y$ 处的 Dirac 测度。
\end{proposition}

\begin{proof}
固定 $(x,t)\in S$。若 $u_r(x,t)=0$，由\eqref{align:ch07-7.89}立即得\eqref{align:ch07-7.139}。设 $y_1\equiv u_r(x,t)>0$。对所有 $y_2>y_1$ 与充分大的 $p\in\mathbb{N}$，置
\begin{align*}
V_p(y)\triangleq p(y-y_1)\chi_{(y_1,y_1+1/p)}(y)+\chi_{[y_1+1/p,y_2)}(y)+p(y_2+1/p-y)\chi_{[y_2,y_2+1/p)}(y),
\end{align*}
\begin{align*}
\Lambda_{V_p}(y)\triangleq \int_0^y V_p'(z)\phi'(z)\,\mathrm{d}z\quad(y\in[0,\infty)).
\end{align*}
则对任意 $y\in[0,\infty)$，
\begin{align}
\lim_{p\to\infty}V_p(y)=\chi_{(y_1,y_2]}(y), \label{align:ch07-7.140a}
\end{align}
\begin{align}
\lim_{p\to\infty}\Lambda_{V_p}(y)=\phi'(y_1)\chi_{(y_1,y_2]}(y)+[\phi'(y_1)-\phi'(y_2)]\chi_{(y_2,\infty)}(y). \label{align:ch07-7.140b}
\end{align}
设 $\{U_p\}$ 为\cref{proposition:ch07-7.5.6}证明中所用序列。由标准逼近论证，可在等式\eqref{align:ch07-7.125}中取 $U=U_p$、$V=V_p$，得
\begin{align*}
\int_{[0,\infty)}[\Lambda_{U_p}(y)-\Lambda_{U_p}^*(x,t)]V_p(y)\,\mathrm{d}\nu_{(x,t)}(y)
=\int_{[0,\infty)}[U_p(y)-U_p^*(x,t)]\Lambda_{V_p}(y)\,\mathrm{d}\nu_{(x,t)}(y).
\end{align*}
由\eqref{align:ch07-7.128}、\eqref{align:ch07-7.132}与\eqref{align:ch07-7.119}，对所有 $y\geqslant0$ 有
\begin{align}
\lim_{p\to\infty}U_p(y)-U_p^*(x,t)=y-y_1, \label{align:ch07-7.141a}
\end{align}
\begin{align}
\lim_{p\to\infty}\Lambda_{U_p}^*(x,t)-\Lambda_{U_p}(y)=\phi(y_1)-\phi(y). \label{align:ch07-7.141b}
\end{align}
由\eqref{align:ch07-7.140a}--\eqref{align:ch07-7.141b}，在 $S$ 中几乎处处得
\begin{align*}
\lim_{p\to\infty}\int_{[0,\infty)}[\Lambda_{U_p}(y)-\Lambda_{U_p}^*]V_p(y)\,\mathrm{d}\nu(y)
=\int_{(y_1,y_2]}[\phi(y)-\phi(y_1)]\,\mathrm{d}\nu(y),
\end{align*}
\begin{align*}
\begin{aligned}
\lim_{p\to\infty}\int_{[0,\infty)}&[U_p(y)-U_p^*]\Lambda_{V_p}(y)\,\mathrm{d}\nu(y)\\
&=\int_{(y_1,y_2]}\phi'(y_1)(y-y_1)\,\mathrm{d}\nu(y)
+[\phi'(y_1)-\phi'(y_2)]\int_{(y_2,\infty)}(y-y_1)\,\mathrm{d}\nu(y),
\end{aligned}
\end{align*}
从而
\begin{align}
\int_{(y_1,y_2]}[\phi(y)-\phi(y_1)-\phi'(y_1)(y-y_1)]\,\mathrm{d}\nu(y)
=[\phi'(y_1)-\phi'(y_2)]\int_{(y_2,\infty)}(y-y_1)\,\mathrm{d}\nu(y). \label{align:ch07-7.142a}
\end{align}
由相同的论证，对任意 $y_0\in(0,y_1)$，
\begin{align}
\begin{aligned}
\int_{(y_0,y_1)}&[\phi(y)-\phi(y_1)-\phi'(y_1)(y-y_1)]\,\mathrm{d}\nu_{(x,t)}(y)\\
&=[\phi'(y_1)-\phi'(y_0)]\int_{[0,y_0]}(y-y_1)\,\mathrm{d}\nu_{(x,t)}(y).
\end{aligned} 
\label{align:ch07-7.142b}
\end{align}
由假设 (A2)，$\phi'$ 至少在区间 $[y_0,y_1]$ 与 $[y_1,y_2]$ 之一上严格单调。设在 $[y_1,y_2]$ 上，另一情形类似。则由\eqref{align:ch07-7.142a}，对几乎所有 $(x,t)\in S$ 有
\begin{align*}
\int_{(y_1,y_2]}|\phi(y)-\phi(y_1)-\phi'(y_1)(y-y_1)|\,\mathrm{d}\nu_{(x,t)}(y)
+|\phi'(y_1)-\phi'(y_2)|\int_{(y_2,\infty)}|y_1-y|\,\mathrm{d}\nu_{(x,t)}(y)=0,
\end{align*}
故 $\operatorname{supp}\nu_{(x,t)}\subseteq[0,y_1]$。因 $\nu_{(x,t)}$ 为概率测度，
\begin{align*}
u_r(x,t)=\int_{[0,u_r(x,t)]}y\,\mathrm{d}\nu_{(x,t)}(y)
=\int_{[0,u_r(x,t)]}[y-u_r(x,t)]\,\mathrm{d}\nu_{(x,t)}(y)+u_r(x,t),
\end{align*}
从而
\begin{align*}
\int_{[0,u_r(x,t)]}|y-u_r(x,t)|\,\mathrm{d}\nu_{(x,t)}(y)=0.
\end{align*}
由上式得 $\operatorname{supp}\nu_{(x,t)}=\{u_r(x,t)\}$，因 $\nu_{(x,t)}$ 为概率测度，\eqref{align:ch07-7.139}得证。
\end{proof}

% SUPPLEMENT: parent=proposition:ch07-7.5.7; kind=proof; location=命题 7.5.7 原证明中的三处省略; basis=原书以“由标准逼近论证可在 (7.125) 中取 $U=U_p$、$V=V_p$”与“$u_r=0$ 时立即得”带过两处、又直接以 (7.140)、(7.141) 的逐点极限推出两个积分等式的极限并取绝对值；此处补出 $V_p$ 的磨光与 (7.140) 的计算、控制收敛的优函数、符号分析，以及 $\operatorname{supp}\nu\subseteq[0,y_1]$ 的两情形论证
\begin{proof}
补出原证明中四处未展开的环节。

第一步：$u_r(x,t)=0$ 的情形。由\eqref{align:ch07-7.89}，$\int_{[0,\infty)}y\,\mathrm{d}\nu_{(x,t)}(y)=0$。因 $\nu_{(x,t)}$ 是 $[0,\infty)$ 上的非负测度且 $y\geqslant0$，得 $\nu_{(x,t)}((0,\infty))=0$，故 $\nu_{(x,t)}=\nu_{(x,t)}(\{0\})\delta_0$；又 $\nu_{(x,t)}$ 为概率测度，故 $\nu_{(x,t)}=\delta_0=\delta_{u_r(x,t)}$。

第二步：$V_p$ 的合法性与\eqref{align:ch07-7.140a}--\eqref{align:ch07-7.140b}。$V_p$ 连续、分段线性且 $0\leqslant V_p\leqslant1$，在 $y_1,y_1+1/p,y_2,y_2+1/p$ 处有折点，其导数 $V_p'=p\chi_{(y_1,y_1+1/p)}-p\chi_{(y_2,y_2+1/p)}$ 分段常值。取标准磨光核 $\theta_\delta$ 与 $V_{p,\delta}\triangleq V_p*\theta_\delta\in C^\infty_c([0,\infty))$（在端点处用截断），则 $V_{p,\delta}\to V_p$、$\Lambda_{V_{p,\delta}}\to\Lambda_{V_p}$ 在 $[0,\infty)$ 上局部一致（后者因 $|\phi'|\leqslant M$）。先在\eqref{align:ch07-7.125}中取 $V=V_{p,\delta}$，再令 $\delta\downarrow0$，由控制收敛（优函数 $M(y+|y_2|+1)$ 属于 $L^1(\nu_{(x,t)})$，见\eqref{align:ch07-7.95}）即得所需等式。关于\eqref{align:ch07-7.140a}：对 $y\in(y_1,y_1+1/p)$ 有 $V_p(y)=p(y-y_1)\to1$ 当 $p\to\infty$ 且 $y>y_1$ 固定；对 $y\in[y_1+1/p,y_2)$ 有 $V_p(y)=1$；对 $y>y_2$ 与 $p$ 充分大（$p>1/(y-y_2)$）有 $V_p(y)=0$；对 $y\leqslant y_1$ 有 $V_p(y)=0$。合并即得 $\lim_pV_p=\chi_{(y_1,y_2]}$。关于\eqref{align:ch07-7.140b}：$\Lambda_{V_p}(y)=\int_0^yV_p'\phi'$，而 $V_p'$ 在 $(y_1,y_1+1/p)$ 上等于 $p$、在 $(y_2,y_2+1/p)$ 上等于 $-p$，这两族核分别逼近 $\delta_{y_1}$ 与 $\delta_{y_2}$；由 $\phi'$ 连续得
\begin{align*}
\Lambda_{V_p}(y)\to\phi'(y_1)\chi_{(y_1,\infty)}(y)-\phi'(y_2)\chi_{(y_2,\infty)}(y)
=\phi'(y_1)\chi_{(y_1,y_2]}(y)+[\phi'(y_1)-\phi'(y_2)]\chi_{(y_2,\infty)}(y).
\end{align*}

第三步：两个积分等式的极限与符号。由\eqref{align:ch07-7.127}得 $|U_p(y)-U_p^*|\leqslant|y|+|U_p^*|$ 且 $|U_p^*|\leqslant\int y\,\mathrm{d}\nu=u_b(x,t)$，又 $|\Lambda_{U_p}(y)-\Lambda_{U_p}^*|\leqslant My+M\int z\,\mathrm{d}\nu(z)=My+Mu_b(x,t)$；又由 $|V_p|\leqslant1$、$|\Lambda_{V_p}|\leqslant M(y+|y_2|)$，被积函数被 $(1+M)(y+u_b(x,t)+|y_2|)\in L^1(\nu_{(x,t)})$ 控制，故控制收敛定理可用；逐点极限\eqref{align:ch07-7.140a}--\eqref{align:ch07-7.141b}可移入积分。需指出\eqref{align:ch07-7.141b}给出的是 $\Lambda_{U_p}^*-\Lambda_{U_p}(y)\to\phi(y_1)-\phi(y)$，故 $\Lambda_{U_p}-\Lambda_{U_p}^*\to\phi(y)-\phi(y_1)$，代入即得第一式的极限 $\int_{(y_1,y_2]}[\phi(y)-\phi(y_1)]\,\mathrm{d}\nu$。

第四步：由\eqref{align:ch07-7.142a}到 $\operatorname{supp}\nu\subseteq[0,y_1]$。设 $\phi'$ 在 $[y_1,y_2]$ 上严格单调。若 $\phi'$ 严格递增，则对 $y\in(y_1,y_2]$ 有 $\phi'(y)>\phi'(y_1)$，故
\begin{align*}
\phi(y)-\phi(y_1)-\phi'(y_1)(y-y_1)=\int_{y_1}^y[\phi'(s)-\phi'(y_1)]\,\mathrm{d}s>0,
\end{align*}
而 $\phi'(y_1)-\phi'(y_2)<0$、$\int_{(y_2,\infty)}(y-y_1)\,\mathrm{d}\nu\geqslant0$。于是\eqref{align:ch07-7.142a}的左端 $\geqslant0$、右端 $\leqslant0$，故两者均为 $0$；特别地
\begin{align*}
[\phi'(y_1)-\phi'(y_2)]\int_{(y_2,\infty)}(y-y_1)\,\mathrm{d}\nu=0,
\end{align*}
而 $\phi'(y_1)\neq\phi'(y_2)$，故 $\nu((y_2,\infty))=0$。若 $\phi'$ 严格递减，则左端 $\leqslant0$、右端 $\geqslant0$，同样得 $\nu((y_2,\infty))=0$。因 $y_2>y_1$ 任意，得 $\nu_{(x,t)}((y_1,\infty))=0$，即 $\operatorname{supp}\nu_{(x,t)}\subseteq[0,y_1]$。

第五步：$\phi'$ 在 $[y_0,y_1]$ 上严格单调的情形。此时由\eqref{align:ch07-7.142b}作同样的符号分析：对 $y\in(y_0,y_1)$ 与严格递增的 $\phi'$，有 $\phi(y)-\phi(y_1)-\phi'(y_1)(y-y_1)=\int_{y_1}^y[\phi'(s)-\phi'(y_1)]\,\mathrm{d}s>0$，而 $\phi'(y_1)-\phi'(y_0)>0$、$\int_{[0,y_0]}(y-y_1)\,\mathrm{d}\nu\leqslant0$，故\eqref{align:ch07-7.142b}两端均为 $0$，于是 $\nu([0,y_0])=0$；由 $y_0\in(0,y_1)$ 任意得 $\operatorname{supp}\nu_{(x,t)}\subseteq[y_1,\infty)$。结合\eqref{align:ch07-7.89}，
\begin{align*}
0=\int_{[0,\infty)}(y-y_1)\,\mathrm{d}\nu_{(x,t)}(y),
\end{align*}
而右端被积函数在 $[y_1,\infty)$ 上非负，故 $\nu_{(x,t)}((y_1,\infty))=0$，即 $\nu_{(x,t)}=\delta_{y_1}$。

第六步：收尾。在第四步的情形下 $\operatorname{supp}\nu_{(x,t)}\subseteq[0,y_1]=[0,u_r(x,t)]$，故 $\nu_{(x,t)}([0,u_r(x,t)])=1$，原证明中的等式 $u_r=\int_{[0,u_r]}(y-u_r)\,\mathrm{d}\nu+u_r$ 成立；于是 $\int_{[0,u_r]}|y-u_r|\,\mathrm{d}\nu=0$，即 $\nu_{(x,t)}$ 集中在其重心 $u_r(x,t)$ 上，得\eqref{align:ch07-7.139}。
\end{proof}

% SUPPLEMENT: parent=proposition:ch07-7.5.7; kind=motivation; location=命题 7.5.7 后的注记; basis=命题 7.5.7 与 (7.139)、(7.142a)--(7.142b) + 为何 (A2) 的严格单调性是关键、$y_1=u_r$ 充当分界点的原因、$V_p$ 的构造思路、该命题只给出几乎处处的结论、以及 $\phi'$ 在某区间为常数时结论失效的边界（注 7.5.3）
\begin{remark}\phantomsection\label{remark:ch07-7.5.7动机}
\begin{enumerate}[(i)]
\item 引入动机。\cref{proposition:ch07-7.5.6}证明了 $\phi(u_r)=\phi^*$，即 $\phi$ 在 $\nu_{(x,t)}$ 上的平均值等于 $\phi$ 在重心处的值。若 $\phi$ 在 $\operatorname{supp}\nu_{(x,t)}$ 上严格凸或严格凹，Jensen 不等式立即迫使 $\nu_{(x,t)}$ 是 Dirac 测度；\cref{proposition:ch07-7.5.7}用 (A2) 提供的局部严格单调性代替凸性，得到同一结论。有了\eqref{align:ch07-7.139}，Young 测度熵解不等式\eqref{align:ch07-7.5}就退化为\cref{definition:ch07-Qtau熵解}中的熵不等式\eqref{align:ch07-7.8}，这正是\cref{theorem:ch07-7.3.2}的内容。
\item (A2) 的作用。\eqref{align:ch07-7.142a}--\eqref{align:ch07-7.142b}是两条关于 $\nu_{(x,t)}$ 的矩恒等式，其左端的被积函数在 $\phi'$ 严格单调时具有确定的符号，右端的系数 $\phi'(y_1)-\phi'(y_2)$（或 $\phi'(y_1)-\phi'(y_0)$）与另一因子同号，于是“同号两端相等”迫使两者为零。若 $\phi'$ 在 $[y_1,y_2]$ 上为常数，则左端恒为零而 $\phi'(y_1)-\phi'(y_2)=0$，\eqref{align:ch07-7.142a}不携带任何信息，(A2) 恰是排除这种退化情形的条件。需指出 (A2) 只要求 $\phi'$ 在 $u_r$ 的某一侧（$[y_1-a,y_1+b]$ 中的某个区间）严格单调，这已足够：$u_r$ 是重心，只需排除质量分布在重心一侧的可能性。
\item $y_1=u_r$ 的含义。取 $y_1$ 为重心是关键：\eqref{align:ch07-7.142a}把 $(y_1,y_2]$ 上的“二阶矩”与 $(y_2,\infty)$ 上的一阶矩联系起来，而 $\operatorname{supp}\nu\subseteq[0,y_1]$ 一旦成立，$\nu$ 的支撑落在重心左侧、重心又等于 $\int y\,\mathrm{d}\nu$，两者只能同时发生当 $\nu$ 集中在 $y_1$。
\item $V_p$ 的构造。$V_p$ 用两个陡峭的斜坡 $p(y-y_1)$ 与 $p(y_2+1/p-y)$ 分别逼近 $\chi_{(y_1,\infty)}$ 与 $\chi_{(y_2,\infty)}$ 并相减，故 $V_p\to\chi_{(y_1,y_2]}$；其导数为 $p\chi_{(y_1,y_1+1/p)}-p\chi_{(y_2,y_2+1/p)}$，于是 $\Lambda_{V_p}$ 的极限\eqref{align:ch07-7.140b}把 $\phi'(y_1)$ 与 $\phi'(y_2)$ 的信息带出来。$y_2>y_1$ 任意是必要的：只有对一切 $y_2>y_1$ 排除 $(y_2,\infty)$ 上的质量，才能得到 $\operatorname{supp}\nu\subseteq[0,y_1]$。
\item 边界。结论只在 $S$ 中几乎处处成立：证明中固定 $(x,t)$ 后所用的等式\eqref{align:ch07-7.125}与\eqref{align:ch07-7.119}本身只在几乎处处意义下成立。若 $\phi'$ 在 $[u_r-a,u_r+b]$（$a,b>0$）上为常数，则\eqref{align:ch07-7.139}可以不成立：此时 $\operatorname{supp}\nu_{(x,t)}$ 可包含 $\phi'=\phi'(u_r)$ 的最大区间（见\cref{remark:ch07-7.5.3}），\cref{theorem:ch07-7.3.2}所要求的熵解性质随之失效，这也是 (A2) 不可去掉的原因。
\end{enumerate}
\end{remark}

\begin{remark}\phantomsection\label{remark:ch07-7.5.2}
由\eqref{align:ch07-7.139}与\cref{proposition:ch05-5.4.1}，可能抽取子列（不重新标号），有
\begin{align}
u_{n_k}\to u_r\quad\text{在 }S\text{ 中几乎处处}. \label{align:ch07-7.143}
\end{align}
特别地，若 $\phi$ 有界，则由控制收敛定理，对任意 $L>0$ 有 $\phi(u_{n_k})\to\phi(u_r)$ 在 $L^1(R_L)$ 中。
\end{remark}

% SUPPLEMENT: parent=remark:ch07-7.5.2; kind=proof; location=注 7.5.2 中“由 (7.139) 与命题 5.4.1”一句; basis=原书未写出把 (7.139) 翻译为“与 $u_r$ 相伴的 Young 测度”这一中间步骤、命题 5.4.1 的哪一款被使用、以及有界 $\phi$ 时 $L^1(R_L)$ 收敛的控制函数
\begin{proof}
补出\eqref{align:ch07-7.143}的两步推导。

第一步：识别 Young 测度。由\eqref{align:ch07-7.139}，$\nu_{(x,t)}=\delta_{u_r(x,t)}$ 对几乎所有 $(x,t)\in S$ 成立。另一方面，与函数 $u_r\in L^1(S)$ 相伴的 Young 测度 $\nu_{u_r}$ 由
\begin{align*}
\nu_{u_r}(\varphi)=\iint_S\varphi(x,t,u_r(x,t))\,\mathrm{d}x\,\mathrm{d}t\quad(\varphi\in\mathcal{C}_0(S\times\mathbb{R}))
\end{align*}
确定（见\cref{definition:ch02-Young测度}），故其离散分解为 $\{\delta_{u_r(x,t)}\}_{(x,t)\in S}$。于是 $\nu=\nu_{u_r}$。

第二步：应用\cref{proposition:ch05-几乎处处收敛}(ii)。取 $X=S$、$\mu$ 为 $S$ 上的 Lebesgue 测度、$Y=\mathbb{R}$、$u=u_r$。由\cref{proposition:ch05-5.4.1}(ii)，若与 $\{u_{n_k}\}$ 相伴的 Young 测度序列窄收敛于 $\nu_{u_r}$，则存在子列使 $u_{n_k}\to u_r$ 于 $S$ 中几乎处处成立。而窄收敛性已由\cref{theorem:ch07-7.5.1}(ii)给出，$\nu=\nu_{u_r}$ 由第一步给出，故\eqref{align:ch07-7.143}成立。

第三步：$L^1(R_L)$ 中的收敛。设 $\phi$ 有界，$\|\phi\|_{L^\infty([0,\infty))}\leqslant\|\phi\|_\infty$。由\eqref{align:ch07-7.143}得 $\phi(u_{n_k})\to\phi(u_r)$ 于 $R_L$ 中几乎处处（$\phi$ 连续）。又
\begin{align*}
|\phi(u_{n_k})-\phi(u_r)|\leqslant2\|\phi\|_\infty\quad\text{于 }R_L,
\end{align*}
且 $\lambda(R_L)=2LT<\infty$，故由控制收敛定理得 $\|\phi(u_{n_k})-\phi(u_r)\|_{L^1(R_L)}\to0$。
\end{proof}

\begin{remark}\phantomsection\label{remark:ch07-7.5.3}
设 (A0)--(A1) 成立，且对某 $a>0$、$b>0$ 与 $(x,t)\in S$，$\phi'$ 在区间 $I_{a,b}=[u_r(x,t)-a,u_r(x,t)+b]$ 上为常数。此情形可证 $\operatorname{supp}\nu_{(x,t)}\subseteq I_{(x,t)}$，其中 $I_{(x,t)}$ 为 $\phi'=\phi'(u_r)$ 的最大区间（见文献~\cite{ref18} 命题 5.9）。
\end{remark}

\begin{proof}[\cref{theorem:ch07-7.3.2}的证明]
由\cref{theorem:ch07-7.3.1}，$u$ 是 (P) 的 Young 测度熵解，故不等式\eqref{align:ch07-7.5}成立。由\eqref{align:ch07-7.5}与\eqref{align:ch07-7.139}得\eqref{align:ch07-7.8}，结论得证。
\end{proof}

% SUPPLEMENT: parent=theorem:ch07-7.3.2; kind=proof; location=定理 7.3.2 原证明后的补证; basis=原书以“由 (7.5) 与 (7.139) 得 (7.8)”一句给出，未写出把 $\nu_{(x,t)}=\delta_{u_r(x,t)}$ 代入 (7.5) 的三项（$E_l^*$、$F_l^*$ 与 $\langle u_s,\cdot\rangle$）的具体计算
\begin{proof}
补出由\eqref{align:ch07-7.5}到\eqref{align:ch07-7.8}的计算。由\eqref{align:ch07-7.139}，对几乎所有 $(x,t)\in S$ 与每个 $l\in[0,\infty)$，
\begin{align*}
E_l^*(x,t)=\int_{[0,\infty)}E_l(y)\,\mathrm{d}\nu_{(x,t)}(y)=E_l(u_r(x,t)),\qquad
F_l^*(x,t)=\int_{[0,\infty)}F_l(y)\,\mathrm{d}\nu_{(x,t)}(y)=F_l(u_r(x,t)),
\end{align*}
故\eqref{align:ch07-7.5}的第一项化为
\begin{align*}
\iint_S[E_l(u_r)\partial_t\zeta+F_l(u_r)\partial_x\zeta]\,\mathrm{d}x\,\mathrm{d}t.
\end{align*}
第二项 $\int_0^T\langle u_s(\cdot,t),\partial_t\zeta(\cdot,t)\rangle\,\mathrm{d}t$ 在\eqref{align:ch07-7.5}与\eqref{align:ch07-7.8}中相同；右端的边界项 $\int_{\mathbb{R}}E_l(u_{0r})\zeta(x,0)\,\mathrm{d}x+\langle u_{0s},\zeta(\cdot,0)\rangle$ 也相同（取 $\Omega=\mathbb{R}$、$\tau=T$）。于是\eqref{align:ch07-7.5}恰为\eqref{align:ch07-7.8}，即 $u$ 是 (P) 的熵解。
\end{proof}

最后证明\cref{lemma:ch07-7.5.4}。

\begin{proof}[\cref{lemma:ch07-7.5.4}的证明]
由标准逼近论证，只需证等式\eqref{align:ch07-7.106a}。

(i) 先在 $U\in C^2([0,\infty))$、$U'\in L^\infty(0,\infty)$ 满足\eqref{align:ch07-7.58}与\eqref{align:ch07-7.91}的假设下证\eqref{align:ch07-7.106a}。设 $\rho\in C^2_c(\mathbb{R})$，$h\in C_c(0,T)$，$\operatorname{supp}\rho\subseteq I_L$（某 $L>0$）。设 $U_{n_k,\rho}$ 由\eqref{align:ch07-7.60}定义，即
\begin{align*}
U_{n,\rho}(t)\triangleq \int_{\mathbb{R}}U(u_n(x,t))\rho(x)\,\mathrm{d}x.
\end{align*}
由\eqref{align:ch07-7.104}得
\begin{align}
\lim_{k\to\infty}\int_0^T U_{n_k,\rho}(t)h(t)\,\mathrm{d}t
=\int_0^T U_\rho^*(t)h(t)\,\mathrm{d}t+M_U\int_0^T\langle\sigma(\cdot,t),\rho\rangle h(t)\,\mathrm{d}t. \label{align:ch07-7.144}
\end{align}
另一方面，由\cref{lemma:ch07-7.4.6}，序列 $\{U_{n_k,\rho}\}$ 在 $BV(0,T)$ 中有界，故存在子列（不重新标号）在 $L^1(0,T)$ 中收敛于 $U_\rho^*+M_U\langle\sigma(\cdot,\cdot),\rho\rangle$（见\eqref{align:ch07-7.144}）。因 $\{U(u_{n_k})\}$ 在 $L^\infty(0,T;L^1(I_L))$ 中有界且 $U^*\in L^\infty(0,T;L^1(I_L))$（见\cref{proposition:ch07-7.5.3}与不等式\eqref{align:ch07-7.101}），条件 $\rho\in C^2_c(\mathbb{R})$ 可放宽为 $\rho\in C_c(\mathbb{R})$。故此情形下\eqref{align:ch07-7.106a}成立。

(ii) 现对所有 $U\in C_b([0,\infty))$ 证\eqref{align:ch07-7.106a}（此时 $M_U=0$）。置
\begin{align*}
U_p:[0,\infty)\to\mathbb{R},\quad U_p(y)\triangleq ((U\chi_{[0,p]})*\theta_p)(y)\quad(y\geqslant0),
\end{align*}
其中 $\theta_p(y)\triangleq p\theta(py)$，$\theta\in C^\infty_c(\mathbb{R})$ 为标准磨光核（$p\in\mathbb{N}$）。则 $\{U_p\}\subseteq C^2_c([0,\infty))$，$U_p\to U$ 在 $[0,\infty)$ 的紧子集上一致收敛，且 $\|U_p\|_{L^\infty(\mathbb{R})}\leqslant\|U\|_{L^\infty(\mathbb{R})}$。此外 $U_p'\in L^\infty(0,\infty)$，$U_p$ 满足\eqref{align:ch07-7.58}与\eqref{align:ch07-7.91}（$M_{U_p}=0$），故由 (i) 得
\begin{align}
\lim_{k\to\infty}\int_0^T\left|\int_{\mathbb{R}}U_p(u_{n_k}(x,t))\rho(x)\,\mathrm{d}x-U_{p,\rho}^*(t)\right|\mathrm{d}t=0, \label{align:ch07-7.145}
\end{align}
其中 $U_p^*\triangleq \int_{[0,\infty)}U_p(y)\,\mathrm{d}\nu(y)$，$U_{p,\rho}^*\triangleq \int_{\mathbb{R}}U_p^*(x,t)\rho(x)\,\mathrm{d}x$（对所有 $\rho\in C_c(\mathbb{R})$、$p\in\mathbb{N}$）。现注意
\begin{align}
\begin{aligned}
\int_0^T&\left|\int_{\mathbb{R}}U(u_{n_k}(x,t))\rho(x)\,\mathrm{d}x-U_\rho^*(t)\right|\mathrm{d}t\\
&\leqslant\iint_S|U(u_{n_k}(x,t))-U_p(u_{n_k}(x,t))|\,|\rho(x)|\,\mathrm{d}x\,\mathrm{d}t\\
&\quad+\int_0^T\left|\int_{\mathbb{R}}U_p(u_{n_k}(x,t))\rho(x)\,\mathrm{d}x-U_{p,\rho}^*(t)\right|\mathrm{d}t
+\int_0^T|U_\rho^*(t)-U_{p,\rho}^*(t)|\,\mathrm{d}t.
\end{aligned} 
\label{align:ch07-7.146}
\end{align}
关于\eqref{align:ch07-7.146}右端第一项，对所有 $p\in\mathbb{N}$ 与 $M>0$，由 Chebyshev 不等式与\eqref{align:ch07-7.71}中第二个不等式得
\begin{align}
\begin{aligned}
\iint_S&|U(u_{n_k})-U_p(u_{n_k})|\,|\rho|\,\mathrm{d}x\,\mathrm{d}t\\
&\leqslant\iint_{\{u_{n_k}\leqslant M\}}|U(u_{n_k})-U_p(u_{n_k})|\,|\rho|\,\mathrm{d}x\,\mathrm{d}t
+\iint_{\{u_{n_k}>M\}}|U(u_{n_k})-U_p(u_{n_k})|\,|\rho|\,\mathrm{d}x\,\mathrm{d}t\\
&\leqslant\|\rho\|_{L^\infty(\mathbb{R})}T\left\{\lambda(\operatorname{supp}\rho)\|U-U_p\|_{L^\infty(0,M)}+\frac{2}{M}\|u_0\|_{\mathscr{R}_f(\mathbb{R})}\|U\|_{L^\infty(\mathbb{R})}\right\}.
\end{aligned} 
\label{align:ch07-7.147}
\end{align}
关于\eqref{align:ch07-7.146}右端第三项，
\begin{align}
\begin{aligned}
\int_0^T|U_\rho^*(t)-U_{p,\rho}^*(t)|\,\mathrm{d}t
&\leqslant\iint_S|U^*-U_p^*|\,|\rho|\,\mathrm{d}x\,\mathrm{d}t\\
&\leqslant\|\rho\|_{L^\infty(\mathbb{R})}\left\{\lambda(\operatorname{supp}\rho)T\|U-U_p\|_{L^\infty(0,M)}
+2\|U\|_{L^\infty(\mathbb{R})}\iint_{\operatorname{supp}\rho\times(0,T)}\nu_{(x,t)}(\{y>M\})\,\mathrm{d}x\,\mathrm{d}t\right\}.
\end{aligned} 
\label{align:ch07-7.148}
\end{align}
由\eqref{align:ch07-7.145}--\eqref{align:ch07-7.148}，对所有 $\rho\in C_c(\mathbb{R})$、$p\in\mathbb{N}$ 与 $M>0$ 得
\begin{align*}
\begin{aligned}
\limsup_{k\to\infty}\int_0^T\mathrm{d}t\,&\left|\int_{\mathbb{R}}U(u_{n_k}(x,t))\rho(x)\,\mathrm{d}x-U_\rho^*(t)\right|\\
&\leqslant2\|\rho\|_{L^\infty(\mathbb{R})}\Big\{\lambda(\operatorname{supp}\rho)T\|U-U_p\|_{L^\infty(0,M)}
+\|U\|_{L^\infty(\mathbb{R})}\left[\frac{T}{M}+\iint_{\operatorname{supp}\rho\times(0,T)}\nu_{(x,t)}(\{y>M\})\,\mathrm{d}x\,\mathrm{d}t\right]\Big\}.
\end{aligned}
\end{align*}
因 $U_p\to U$ 在 $[0,\infty)$ 的紧子集上一致收敛，令 $p\to\infty$ 得
\begin{align}
\limsup_{k\to\infty}\int_0^T\left|\int_{\mathbb{R}}U(u_{n_k})\rho(x)\,\mathrm{d}x-U_\rho^*(t)\right|\mathrm{d}t
\leqslant2\|\rho\|_{L^\infty(\mathbb{R})}\|U\|_{L^\infty(\mathbb{R})}\left\{\frac{T}{M}+\iint_{\operatorname{supp}\rho\times(0,T)}\nu_{(x,t)}(\{y>M\})\,\mathrm{d}x\,\mathrm{d}t\right\}. \label{align:ch07-7.149}
\end{align}
因 $\nu_{(x,t)}$ 为概率测度，由控制收敛定理得
\begin{align*}
\lim_{M\to\infty}\iint_{\operatorname{supp}\rho\times(0,T)}\nu_{(x,t)}(\{y>M\})\,\mathrm{d}x\,\mathrm{d}t=0.
\end{align*}
在\eqref{align:ch07-7.149}中令 $M\to\infty$，即得\eqref{align:ch07-7.106a}。

(iii) 现设 $U\in C([0,\infty))$ 为任意满足\eqref{align:ch07-7.91}的函数。如\cref{proposition:ch07-7.5.3}证明中那样，设 $f_{j1},f_{j2}\in C^2([0,\infty))$ 满足\eqref{align:ch07-7.100}（$j\in\mathbb{N}$，$j\geqslant[y_\varepsilon]+1$）。则
\begin{align}
U=U_{j1}+U_{j2},\quad U^*=U_{j1}^*+U_{j2}^*,\quad U_\rho^*=U_{j1,\rho}^*+U_{j2,\rho}^*, \label{align:ch07-7.150}
\end{align}
其中 $U_{jl}\triangleq f_{jl}U$，$U_{jl}^*\triangleq \int_{[0,\infty)}U_{jl}(y)\,\mathrm{d}\nu(y)$，$U_{jl,\rho}^*\triangleq \int_{\mathbb{R}}U_{jl}^*(x,t)\rho(x)\,\mathrm{d}x$（$\rho\in C_c(\mathbb{R})$，$l=1,2$）。因此
\begin{align}
\begin{aligned}
\int_0^T&\left|\int_{\mathbb{R}}U(u_{n_k}(x,t))\rho(x)\,\mathrm{d}x-U_\rho^*(t)-M_U\langle\sigma(\cdot,t),\rho\rangle\right|\mathrm{d}t\\
&\leqslant\int_0^T\left|\int_{\mathbb{R}}U_{j1}(u_{n_k}(x,t))\rho(x)\,\mathrm{d}x-U_{j1,\rho}^*(t)\right|\mathrm{d}t\\
&\quad+\int_0^T\left|\int_{\mathbb{R}}U_{j2}(u_{n_k}(x,t))\rho(x)\,\mathrm{d}x-U_{j2,\rho}^*(t)-M_U\langle\sigma(\cdot,t),\rho\rangle\right|\mathrm{d}t=:I_{jk1}+I_{jk2}.
\end{aligned} 
\label{align:ch07-7.151}
\end{align}
因 $U_{j1}\in C_b([0,\infty))$ 从而 $M_{U_{j1}}=0$，由 (ii) 与\eqref{align:ch07-7.106a}得
\begin{align}
\lim_{k\to\infty}I_{jk1}=0\quad(j\in\mathbb{N},\ j\geqslant[y_\varepsilon]+1). \label{align:ch07-7.152}
\end{align}
另一方面，
\begin{align}
\begin{aligned}
I_{jk2}\leqslant{}&\int_0^T\left|\int_{\mathbb{R}}[U_{j2}(u_{n_k}(x,t))-M_U f_{j2}(u_{n_k}(x,t))u_{n_k}(x,t)]\rho(x)\,\mathrm{d}x\right|\mathrm{d}t\\
&+M_U\int_0^T\mathrm{d}t\left|\int_{\mathbb{R}}f_{j2}(u_{n_k}(x,t))u_{n_k}(x,t)\rho(x)\,\mathrm{d}x-\int_{\mathbb{R}}F_{j2}^*(x,t)\rho(x)\,\mathrm{d}x-\langle\sigma(\cdot,t),\rho\rangle\right|\\
&+\iint_S|U_{j2}^*-M_U F_{j2}^*|(x,t)\,|\rho(x)|\,\mathrm{d}x\,\mathrm{d}t,
\end{aligned} 
\label{align:ch07-7.153}
\end{align}
其中 $F_{j2}^*\triangleq \int_{[0,\infty)}f_{j2}(y)y\,\mathrm{d}\nu(y)$。

分别处理\eqref{align:ch07-7.153}右端三项。由\eqref{align:ch07-7.99}与\eqref{align:ch07-7.71}中第二个不等式得
\begin{align}
\int_0^T\left|\int_{\mathbb{R}}[U_{j2}(u_{n_k}(x,t))-M_U f_{j2}(u_{n_k}(x,t))u_{n_k}(x,t)]\rho(x)\,\mathrm{d}x\right|\mathrm{d}t
\leqslant\varepsilon\|\rho\|_{L^\infty(\mathbb{R})}T\|u_0\|_{\mathscr{R}_f(\mathbb{R})}. \label{align:ch07-7.154}
\end{align}
此外，映射 $y\mapsto f_{j2}(y)y$ 属于 $C^2([0,\infty))$，导数有界，满足\eqref{align:ch07-7.58}与\eqref{align:ch07-7.91}（$M_U=1$）。故由 (i) 得
\begin{align}
\lim_{k\to\infty}\int_0^T\mathrm{d}t\left|\int_{\mathbb{R}}f_{j2}(u_{n_k}(x,t))u_{n_k}(x,t)\rho(x)\,\mathrm{d}x-\int_{\mathbb{R}}F_{j2}^*(x,t)\rho(x)\,\mathrm{d}x-\langle\sigma(\cdot,t),\rho\rangle\right|=0. \label{align:ch07-7.155}
\end{align}
现证
\begin{align}
\lim_{j\to\infty}\iint_S|U_{j2}^*-M_U F_{j2}^*|(x,t)\,|\rho(x)|\,\mathrm{d}x\,\mathrm{d}t=0. \label{align:ch07-7.156}
\end{align}
对几乎所有 $(x,t)\in S$ 有
\begin{align*}
|U_{j2}^*(x,t)-M_U F_{j2}^*(x,t)|
&\leqslant\int_{[0,\infty)}|U_{j2}(y)-M_U f_{j2}(y)y|\,\mathrm{d}\nu_{(x,t)}(y)\\
&=\int_{[0,\infty)}f_{j2}(y)|U(y)-M_U y|\,\mathrm{d}\nu_{(x,t)}(y)
\leqslant\varepsilon\int_{[j,\infty)}y\,\mathrm{d}\nu_{(x,t)}(y)\leqslant\varepsilon u_b(x,t)
\end{align*}
（见\eqref{align:ch07-7.95}）。因 $u_b\in L^\infty(0,T;L^1(\mathbb{R}))$ 且对几乎所有 $(x,t)\in S$ 有 $\lim_{j\to\infty}\int_{[j,\infty)}y\,\mathrm{d}\nu_{(x,t)}(y)=0$，由控制收敛定理得等式\eqref{align:ch07-7.156}。

由\eqref{align:ch07-7.153}--\eqref{align:ch07-7.156}得
\begin{align*}
\limsup_{j\to\infty}\left(\limsup_{k\to\infty}I_{jk2}\right)\leqslant\varepsilon\|\rho\|_{L^\infty(\mathbb{R})}T\|u_0\|_{\mathscr{R}_f(\mathbb{R})},
\end{align*}
故由\eqref{align:ch07-7.151}--\eqref{align:ch07-7.152}得
\begin{align*}
\limsup_{k\to\infty}\int_0^T\left|\int_{\mathbb{R}}[U(u_{n_k})-U^*](x,t)\rho(x)\,\mathrm{d}x-M_U\langle\sigma(\cdot,t),\rho\rangle\right|\mathrm{d}t
\leqslant\varepsilon T\|u_0\|_{\mathscr{R}_f(\mathbb{R})}\|\rho\|_{L^\infty(\mathbb{R})}.
\end{align*}
由 $\varepsilon$ 任意，上式给出\eqref{align:ch07-7.106a}。证毕。
\end{proof}

% SUPPLEMENT: parent=lemma:ch07-7.5.4; kind=proof; location=引理 7.5.4 原证明中“由标准逼近论证，只需证等式 (7.106a)”与第 (i) 步的极限识别; basis=原书未写出由一般 $\zeta\in C([0,T];C_c(\mathbb{R}))$ 到 $\rho(x)$ 型试验函数的化归、未写出 $BV$ 有界如何给出 $L^1(0,T)$ 中的收敛子列、也未写出该子列的极限为何恰是 $U_\rho^*+M_U\langle\sigma(\cdot,\cdot),\rho\rangle$ 以及 $\rho$ 的正则性可从 $C^2_c$ 放宽到 $C_c$ 的理由
\begin{proof}
补出化归与第 (i) 步的细节。

第一步：化归到 $\rho(x)$ 型试验函数。设\eqref{align:ch07-7.106a}对一切 $\rho\in C_c(\mathbb{R})$ 成立，要证\eqref{align:ch07-7.105}对一切 $\zeta\in C([0,T];C_c(\mathbb{R}))$ 成立。由\eqref{align:ch07-7.101}，存在与 $k$ 无关的常数 $C_1$ 使
\begin{align*}
\left|\int_{\mathbb{R}}U(u_{n_k}(x,t))\zeta(x,t)\,\mathrm{d}x\right|\leqslant C_1\|\zeta\|_{L^\infty(S)},
\end{align*}
又 $|U^*|\leqslant M(u_b+1)$ 给出 $\left|\int_{\mathbb{R}}U^*(x,t)\zeta(x,t)\,\mathrm{d}x\right|\leqslant C_2\|\zeta\|_{L^\infty(S)}$，且 $|\langle\sigma(\cdot,t),\zeta(\cdot,t)\rangle|\leqslant\|\sigma(\cdot,t)\|_{\mathscr{R}_f}\|\zeta\|_{L^\infty(S)}$。故\eqref{align:ch07-7.105}左端的被积函数被 $C\|\zeta\|_{L^\infty(S)}$ 控制，其中 $C$ 与 $k$、$\zeta$ 无关。于是只需对属于 $C([0,T])\otimes C_c(\mathbb{R})$ 的有限和证\eqref{align:ch07-7.105}；由三角不等式可进一步只须对 $\zeta(x,t)=h(t)\rho(x)$（$h\in C([0,T])$，$\rho\in C_c(\mathbb{R})$）证明。此时\eqref{align:ch07-7.105}的左端不超过
\begin{align*}
\int_0^T|h(t)|\left|\int_{\mathbb{R}}U(u_{n_k}(x,t))\rho(x)\,\mathrm{d}x-U_\rho^*(t)-M_U\langle\sigma(\cdot,t),\rho\rangle\right|\mathrm{d}t
\leqslant\|h\|_{L^\infty(0,T)}\int_0^T|\cdots|\,\mathrm{d}t,
\end{align*}
由\eqref{align:ch07-7.106a}即得\eqref{align:ch07-7.105}。

第二步：$BV$ 有界给出 $L^1(0,T)$ 中的收敛子列。由\cref{lemma:ch07-7.4.6}，$\|U_{n_k,\rho}\|_{BV(0,T)}\leqslant C$（$k\in\mathbb{N}$，$C$ 与 $k$ 无关）。由 $BV(0,T)\subseteq L^\infty(0,T)$ 与 $\|U_{n_k,\rho}\|_{L^1(0,T)}$ 的一致有界性（见\eqref{align:ch07-7.65}），$\{U_{n_k,\rho}\}$ 在 $BV(0,T)$ 的单位球中有界；$BV(0,T)$ 紧嵌入 $L^1(0,T)$，故存在子列在 $L^1(0,T)$ 中收敛于某 $g\in L^1(0,T)$。

第三步：识别极限 $g$。对每个 $h\in C_c(0,T)$，\eqref{align:ch07-7.144}给出
\begin{align*}
\int_0^TU_{n_k,\rho}(t)h(t)\,\mathrm{d}t\to\int_0^T\Big\{U_\rho^*(t)+M_U\langle\sigma(\cdot,t),\rho\rangle\Big\}h(t)\,\mathrm{d}t.
\end{align*}
而由第二步的 $L^1(0,T)$ 收敛，同一积分收敛于 $\int_0^Tg(t)h(t)\,\mathrm{d}t$。因 $h\in C_c(0,T)$ 任意，得 $g=U_\rho^*+M_U\langle\sigma(\cdot,\cdot),\rho\rangle$ 于 $(0,T)$ 中几乎处处；特别地 $g\in L^1(0,T)$ 且第二步的收敛就是 $L^1(0,T)$ 中的收敛。

第四步：$\rho\in C^2_c(\mathbb{R})$ 放宽为 $C_c(\mathbb{R})$。设 $\rho\in C_c(\mathbb{R})$，取 $\{\rho_m\}\subseteq C^2_c(\mathbb{R})$ 使 $\operatorname{supp}\rho_m\subseteq I_{L+1}$ 且 $\rho_m\to\rho$ 于 $L^\infty$。对每个 $m$，第三步给出
\begin{align*}
\lim_{k\to\infty}\int_0^T\left|\int_{\mathbb{R}}U(u_{n_k})\rho_m\,\mathrm{d}x-U_{\rho_m}^*(t)-M_U\langle\sigma(\cdot,t),\rho_m\rangle\right|\mathrm{d}t=0.
\end{align*}
由\eqref{align:ch07-7.101}与 $|U^*|\leqslant M(u_b+1)$、$\|\sigma(\cdot,t)\|\leqslant\|\sigma\|$，存在与 $k$、$m$ 无关的常数 $C_3$（依赖 $\|\rho\|_{L^\infty}$ 与 $L$）使被积函数中的差用 $\rho_m-\rho$ 替换后不超过 $C_3\|\rho_m-\rho\|_{L^\infty}$；故先取 $m$ 充分大、再令 $k\to\infty$，即得 $\rho\in C_c(\mathbb{R})$ 情形的\eqref{align:ch07-7.106a}。
\end{proof}

% SUPPLEMENT: parent=lemma:ch07-7.5.4; kind=proof; location=引理 7.5.4 原证明第 (ii)、(iii) 步的定量细节; basis=原书 (7.147)、(7.148) 中的 Chebyshev 估计与 $\lambda(\operatorname{supp}\rho)T$ 的来源未展开，(7.153) 的三项分解与最后的 $\varepsilon$-论证只给结果；此处补出逐项估计与被积函数的控制
\begin{proof}
补出第 (ii)、(iii) 步的定量细节。

第一步：\eqref{align:ch07-7.147}。由 $\|U_p\|_{L^\infty}\leqslant\|U\|_{L^\infty}$ 得 $|U(u_{n_k})-U_p(u_{n_k})|\leqslant2\|U\|_{L^\infty(\mathbb{R})}$。故
\begin{align*}
\iint_{\{u_{n_k}>M\}}|U(u_{n_k})-U_p(u_{n_k})|\,|\rho|\,\mathrm{d}x\,\mathrm{d}t
\leqslant2\|U\|_{L^\infty(\mathbb{R})}\|\rho\|_{L^\infty(\mathbb{R})}\iint_{\{u_{n_k}>M\}}\!\!\mathrm{d}x\,\mathrm{d}t,
\end{align*}
而由 Chebyshev 不等式与\eqref{align:ch07-7.71}中第二个不等式，
\begin{align*}
\iint_{\{u_{n_k}>M\}}\mathrm{d}x\,\mathrm{d}t\leqslant\frac1M\iint_Su_{n_k}\,\mathrm{d}x\,\mathrm{d}t\leqslant\frac{T}{M}\|u_0\|_{\mathscr{R}_f(\mathbb{R})}.
\end{align*}
在 $\{u_{n_k}\leqslant M\}$ 上，$|U(u_{n_k})-U_p(u_{n_k})|\leqslant\|U-U_p\|_{L^\infty(0,M)}$，且 $\iint_S|\rho|\,\mathrm{d}x\,\mathrm{d}t=\lambda(\operatorname{supp}\rho)T\|\rho\|_{L^\infty(\mathbb{R})}$（因 $\operatorname{supp}\rho$ 紧且 $t$ 区间长为 $T$）。合并即得\eqref{align:ch07-7.147}。

第二步：\eqref{align:ch07-7.148}。由\eqref{align:ch07-7.106b}与\eqref{align:ch07-7.93}，
\begin{align*}
|U_\rho^*(t)-U_{p,\rho}^*(t)|\leqslant\int_{\mathbb{R}}|U^*(x,t)-U_p^*(x,t)|\,|\rho(x)|\,\mathrm{d}x
\leqslant\int_{\mathbb{R}}|\rho(x)|\int_{[0,\infty)}|U_p(y)-U(y)|\,\mathrm{d}\nu_{(x,t)}(y)\,\mathrm{d}x.
\end{align*}
把 $y$ 积分拆为 $\{0\leqslant y\leqslant M\}$ 与 $\{y>M\}$ 两部分：前者被 $\|U-U_p\|_{L^\infty(0,M)}$ 控制，后者被 $2\|U\|_{L^\infty(\mathbb{R})}\nu_{(x,t)}(\{y>M\})$ 控制（因 $\nu_{(x,t)}$ 为概率测度）。对 $t$ 积分并用 $\int_0^T\mathrm{d}t\int_{\mathbb{R}}\mathrm{d}x=\iint_S$ 即得\eqref{align:ch07-7.148}。

第三步：\eqref{align:ch07-7.153} 的分解。由\eqref{align:ch07-7.150}，
\begin{align*}
&\int_{\mathbb{R}}U_{j2}(u_{n_k})\rho\,\mathrm{d}x-U_{j2,\rho}^*(t)-M_U\langle\sigma(\cdot,t),\rho\rangle\\
&=\int_{\mathbb{R}}[U_{j2}(u_{n_k})-M_Uf_{j2}(u_{n_k})u_{n_k}]\rho\,\mathrm{d}x
+M_U\Big\{\int_{\mathbb{R}}f_{j2}(u_{n_k})u_{n_k}\rho\,\mathrm{d}x-\int_{\mathbb{R}}F_{j2}^*\rho\,\mathrm{d}x-\langle\sigma(\cdot,t),\rho\rangle\Big\}\\
&\quad-M_U\int_{\mathbb{R}}[F_{j2}^*(x,t)-f_{j2}(u_{n_k}(x,t))u_{n_k}(x,t)]\rho\,\mathrm{d}x
+M_U\int_{\mathbb{R}}F_{j2}^*\rho\,\mathrm{d}x-M_U\int_{\mathbb{R}}F_{j2}^*\rho\,\mathrm{d}x,
\end{align*}
整理（把最后两项合并入第三项）即得\eqref{align:ch07-7.153}。其中用到
\begin{align*}
U_{j2}^*(x,t)-M_UF_{j2}^*(x,t)=\int_{[0,\infty)}[U_{j2}(y)-M_Uf_{j2}(y)y]\,\mathrm{d}\nu_{(x,t)}(y),
\end{align*}
这正是原证明中第三个积分的形式。

第四步：\eqref{align:ch07-7.154} 与最后的 $\varepsilon$-论证。在 $\operatorname{supp}f_{j2}\subseteq[j,\infty)$ 上，由 $j\geqslant[y_\varepsilon]+1$ 得 $y>y_\varepsilon$，故由\eqref{align:ch07-7.99}，
\begin{align*}
|U_{j2}(y)-M_Uf_{j2}(y)y|=f_{j2}(y)|U(y)-M_Uy|\leqslant\varepsilon yf_{j2}(y)\leqslant\varepsilon y.
\end{align*}
对 $y=u_{n_k}(x,t)$ 积分并用\eqref{align:ch07-7.71}中第二个不等式即得\eqref{align:ch07-7.154}。最后，由\eqref{align:ch07-7.151}--\eqref{align:ch07-7.156}，
\begin{align*}
\limsup_{k\to\infty}J_{jk}\leqslant\lim_{k\to\infty}I_{jk1}+\limsup_{k\to\infty}I_{jk2}\leqslant\varepsilon\|\rho\|_{L^\infty(\mathbb{R})}T\|u_0\|_{\mathscr{R}_f(\mathbb{R})},
\end{align*}
其中 $J_{jk}$ 记\eqref{align:ch07-7.151}的左端。上式对每个 $j\geqslant[y_\varepsilon]+1$ 成立，故由\eqref{align:ch07-7.152}先令 $k\to\infty$、再令 $j\to\infty$，再由 $\varepsilon>0$ 的任意性，即得\eqref{align:ch07-7.106a}。
\end{proof}

\end{document}
